\documentclass[12pt]{article}
\usepackage{amsmath}
\usepackage{amsfonts}
\usepackage{float}
\usepackage{fancyhdr}
\usepackage{graphicx}
\usepackage{setspace}
\usepackage[colorlinks=true,linkcolor=blue, citecolor=red]{hyperref}
\usepackage{url}
\usepackage[top=.75in, left=.75in, right=.75in, bottom=1in]{geometry}

% For algorithm
\usepackage{algorithm}
\usepackage{algpseudocode}

% ============ CODE ============
\usepackage{listings}
\usepackage{xcolor}
\definecolor{codegreen}{rgb}{0,0.6,0}
\definecolor{codegray}{rgb}{0.5,0.5,0.5}
\definecolor{codepurple}{rgb}{0.58,0,0.82}
\definecolor{backcolour}{rgb}{0.95,0.95,0.92}

% Styling for the code.
\lstdefinestyle{mystyle}{
    backgroundcolor=\color{backcolour},   
    commentstyle=\color{codegreen},
    keywordstyle=\color{magenta},
    numberstyle=\tiny\color{codegray},
    stringstyle=\color{codepurple},
    basicstyle=\ttfamily\footnotesize,
    breakatwhitespace=false,         
    breaklines=true,                 
    captionpos=b,                    
    keepspaces=true,                 
    numbers=left,                    
    numbersep=5pt,                  
    showspaces=false,                
    showstringspaces=false,
    showtabs=false,                  
    tabsize=2
}
\lstset{style=mystyle}

% Table packages and Unicode / Font support
\usepackage{array}
\usepackage{longtable}
\usepackage{enumitem}
\usepackage{booktabs}
\usepackage[utf8]{inputenc}
\DeclareUnicodeCharacter{2264}{\ensuremath{\le}}
\DeclareUnicodeCharacter{2194}{\ensuremath{\leftrightarrow}}
\usepackage[T1]{fontenc}
\usepackage{lmodern}
\graphicspath{{../../../../templates/}{../../prototype/}{figures/}{./}}
\setlength{\emergencystretch}{3em}
\setlist[enumerate]{label=\arabic*.,leftmargin=*}
\setlist[itemize]{leftmargin=*}

% ============ HEADER AND FOOTER ============
\setlength{\headheight}{29.43912pt}

\pagestyle{fancy}
\fancyfoot{}
\fancyfoot[R]{Page \thepage}

\setlength{\parindent}{0pt}
\renewcommand{\baselinestretch}{1.5}

\newcommand{\coursename}{Introduction of Software Engineering}
\newcommand{\reportname}{Studitopia}
\newcommand{\reporttitle}{\href{https://github.com/RoundFloor/Studitopia.git}{Studitopia}}

\newcommand{\studentname}{%
Nguyen Thanh An (24125024)\\
Nguyen Chi Bao (24125025)\\
Tran Nguyen Khanh Binh (24125026)\\
Tran Canh Anh Tuan (24125048)\\
Nguyen Le Quoc Huy (24125100)
}
\newcommand{\teachername}{Ho Tuan Thanh, Nguyen Thi Minh Tuyen}

% Header
\lhead{\reporttitle}
\rhead{
University of Science - VNUHCM\\
\coursename
}

% Footer
\newcommand{\leftfooter}{}
\lfoot{\leftfooter}

% ============ DOCUMENT ============
\begin{document}

\pagenumbering{roman}
\begin{titlepage}
\newcommand{\HRule}{\rule{\linewidth}{0.5mm}}
\centering

\textsc{\LARGE VNU-HCM}\\[0.5cm]
\textsc{\Large University of Science}\\[0.5cm]
\textsc{\large Faculty of Information Technology}\\[0.5cm]
\textsc{Computer Science}\\[0.5cm]

\HRule \\[0.4cm]
{ 
\huge{\bfseries{\reporttitle}}\\[0.5cm]
\large{\bfseries{Authentication and Account --- Use Case Diagram}}
}\\[0.4cm]
\HRule \\[0.5cm]

\textbf{\large Subject: \coursename}\\[0.5cm]

\begin{minipage}[t]{0.55\textwidth}
\begin{flushleft} \large
\emph{Students:}\\
\studentname
\end{flushleft}
\end{minipage}
~
\begin{minipage}[t]{0.4\textwidth}
\begin{flushright} \large
\emph{Supervising Teachers:} \\
\teachername
\end{flushright}
\end{minipage}\\[1cm]

{\large \today}\\[1cm]

\includegraphics[scale=.20]{img/hcmus-logo.png}\\[1cm] 

\vfill
\end{titlepage}	

\tableofcontents
\pagebreak

\pagenumbering{arabic}
\setcounter{page}{1}

\hypertarget{sec:authentication-and-account-use-case-diagram}{}
\addcontentsline{toc}{section}{Authentication and Account --- Use Case Diagram}
\section*{Authentication and Account --- Use Case Diagram}
\par\medskip

\textbf{Version:} \texttt{1.0}
\par\medskip

\textbf{Performed by:} [Nguyen Thanh An] | \textbf{Edited by:} [Nguyen Thanh An]
\par\medskip

Merged from the focused use case diagrams in this functional group.
\par\medskip

\begin{center}
\includegraphics[width=0.85\textwidth,height=0.75\textheight,keepaspectratio]{"figures/diagram_1.png"}
\end{center}
\par\medskip

\textbf{Diagram notation.} Person shape = human actor; cloud = supporting system; blue = base use case; purple = included; green = extending; orange = specialized; solid line = association; dashed arrow = include or extend; labelled solid arrow = generalization.
\par\medskip

\pagebreak
\hypertarget{sec:uc-01-register-account-and-verify-university-email}{}
\addcontentsline{toc}{section}{UC-01 --- Register Account and Verify University Email}
\section*{UC-01 --- Register Account and Verify University Email}
\par\medskip

\textbf{Version:} \texttt{2.0}
\par\medskip

\textbf{Performed by:} [Nguyen Thanh An] | \textbf{Edited by:} [Nguyen Thanh An]
\par\medskip

\begingroup
\small
\renewcommand{\baselinestretch}{1.0}\selectfont
\renewcommand{\arraystretch}{1.15}
\setlength{\tabcolsep}{6pt}
\begin{longtable}{|>{\raggedright\arraybackslash}p{0.09\textwidth}|>{\raggedright\arraybackslash}p{0.84\textwidth}|}
\hline
\textbf{Field} & \textbf{Value} \\ \hline
\endhead
Use Case ID & UC-01 \\ \hline
Name & Register Account and Verify University Email \\ \hline
Primary Actor & New University Student \\ \hline
Supporting Actors & Email Delivery Service; Google OAuth 2.0 \\ \hline
Feature Group & Authentication \\ \hline
Priority & High \\ \hline
Related User Stories & US-01, US-07 \\ \hline
Related Requirements & FR-01.1 to FR-01.45; FR-07.1 to FR-07.6; NFR-SEC-2, NFR-SEC-3, NFR-SEC-8, NFR-USE-1 \\ \hline
Source Documents & \texttt{proposals/use\_cases\_model/UC-01-register-account.md}; \texttt{proposals/use\_cases\_description.md} (UC-01); \texttt{proposals/use\_cases\_model.md}; \texttt{proposals/use\_case\_diagrams.md}; \texttt{PA3-Group05/use\_case.md} (section 1); \texttt{proposals/overview.md}; \texttt{proposals/functional\_requirement.md}; \texttt{proposals/non\_functional\_requirement.md}; \texttt{proposals/userstory.md}; \texttt{proposals/feature\_flow.md}; \texttt{proposals/screen\_flow.md}; \texttt{proposals/API.md}; \texttt{proposals/response\_code.md}; \texttt{proposals/settings.md} \\ \hline
\end{longtable}
\endgroup
\par\medskip

\hypertarget{sec:part-a-use-case-model}{}
\addcontentsline{toc}{subsection}{Part A --- Use Case Model}
\subsection*{Part A --- Use Case Model}
\par\medskip

\hypertarget{sec:a1-focused-use-case-diagram}{}
\addcontentsline{toc}{subsubsection}{A1. Focused Use Case Diagram}
\subsubsection*{A1. Focused Use Case Diagram}
\par\medskip

\begin{center}
\includegraphics[width=0.85\textwidth,height=0.75\textheight,keepaspectratio]{"figures/UC-01-register-account_1.png"}
\end{center}
\par\medskip

\textbf{Diagram notation.} Person shape = human actor; cloud = supporting system; blue = selected base use case; purple = included; green = extending; orange = specialized; solid line = association; dashed arrow = include or extend; labelled solid arrow = generalization.
\par\medskip

\hypertarget{sec:a2-actor-register}{}
\addcontentsline{toc}{subsubsection}{A2. Actor Register}
\subsubsection*{A2. Actor Register}
\par\medskip

\begingroup
\small
\renewcommand{\baselinestretch}{1.0}\selectfont
\renewcommand{\arraystretch}{1.15}
\setlength{\tabcolsep}{6pt}
\begin{longtable}{|>{\raggedright\arraybackslash}p{0.14\textwidth}|>{\raggedright\arraybackslash}p{0.12\textwidth}|>{\raggedright\arraybackslash}p{0.66\textwidth}|}
\hline
\textbf{Actor} & \textbf{Type} & \textbf{Description} \\ \hline
\endhead
New University Student & Primary human actor & A prospective user who creates an account and verifies a university email address. \\ \hline
University Student & General human actor & Parent actor of New University Student; specialized actors inherit the general student associations. \\ \hline
Email Delivery Service & Supporting system & Delivers time-limited OTP codes for account verification. \\ \hline
Google OAuth 2.0 & Supporting system & Provides the optional external authentication and registration path. \\ \hline
\end{longtable}
\endgroup
\par\medskip

\hypertarget{sec:a3-relationships}{}
\addcontentsline{toc}{subsubsection}{A3. Relationships}
\subsubsection*{A3. Relationships}
\par\medskip

\hypertarget{sec:a3-1-associations}{}
\paragraph*{A3.1 Associations}
\par\medskip

\begingroup
\small
\renewcommand{\baselinestretch}{1.0}\selectfont
\renewcommand{\arraystretch}{1.15}
\setlength{\tabcolsep}{6pt}
\begin{longtable}{|>{\raggedright\arraybackslash}p{0.15\textwidth}|>{\raggedright\arraybackslash}p{0.34\textwidth}|>{\raggedright\arraybackslash}p{0.10\textwidth}|>{\raggedright\arraybackslash}p{0.33\textwidth}|}
\hline
\textbf{Actor} & \textbf{Use Case} & \textbf{Type} & \textbf{Description} \\ \hline
\endhead
New University Student & UC-01 Register Account and Verify University Email & Primary & Starts registration and email verification. \\ \hline
Google OAuth 2.0 & G-02 Register with Google & Supporting & Authenticates the Google registration path. \\ \hline
Email Delivery Service & I-02 Verify Email OTP & Supporting & Delivers the OTP to the university email address. \\ \hline
\end{longtable}
\endgroup
\par\medskip

\hypertarget{sec:a3-2-generalization-and-specialization}{}
\paragraph*{A3.2 Generalization and Specialization}
\par\medskip

\begingroup
\small
\renewcommand{\baselinestretch}{1.0}\selectfont
\renewcommand{\arraystretch}{1.15}
\setlength{\tabcolsep}{6pt}
\begin{longtable}{|>{\raggedright\arraybackslash}p{0.16\textwidth}|>{\raggedright\arraybackslash}p{0.23\textwidth}|>{\raggedright\arraybackslash}p{0.11\textwidth}|>{\raggedright\arraybackslash}p{0.41\textwidth}|}
\hline
\textbf{Child} & \textbf{Parent} & \textbf{Type} & \textbf{Inherited Behavior or Association} \\ \hline
\endhead
New University Student & University Student & Actor generalization & Inherits the general student associations and adds UC-01. \\ \hline
G-01 Register with University Email & UC-01 Register Account and Verify University Email & Use case generalization & Inherits the UC-01 actor association and I-01 include; adds the mandatory OTP include. \\ \hline
G-02 Register with Google & UC-01 Register Account and Verify University Email & Use case generalization & Inherits the UC-01 actor association and I-01 include; uses Google OAuth instead of OTP. \\ \hline
\end{longtable}
\endgroup
\par\medskip

\hypertarget{sec:a3-3-include}{}
\paragraph*{A3.3 Include}
\par\medskip

\begingroup
\small
\renewcommand{\baselinestretch}{1.0}\selectfont
\renewcommand{\arraystretch}{1.15}
\setlength{\tabcolsep}{6pt}
\begin{longtable}{|>{\raggedright\arraybackslash}p{0.26\textwidth}|>{\raggedright\arraybackslash}p{0.16\textwidth}|>{\raggedright\arraybackslash}p{0.50\textwidth}|}
\hline
\textbf{Base Use Case} & \textbf{Included Use Case} & \textbf{Meaning} \\ \hline
\endhead
UC-01 Register Account and Verify University Email & I-01 Validate Registration Data & Validate required fields, username, email, password, confirmation, CAPTCHA, and uniqueness rules. \\ \hline
G-01 Register with University Email & I-02 Verify Email OTP & Send and verify a 6-digit OTP with a 5-minute expiry. \\ \hline
\end{longtable}
\endgroup
\par\medskip

\hypertarget{sec:a3-4-extend}{}
\paragraph*{A3.4 Extend}
\par\medskip

\begingroup
\small
\renewcommand{\baselinestretch}{1.0}\selectfont
\renewcommand{\arraystretch}{1.15}
\setlength{\tabcolsep}{6pt}
\begin{longtable}{|>{\raggedright\arraybackslash}p{0.10\textwidth}|>{\raggedright\arraybackslash}p{0.11\textwidth}|>{\raggedright\arraybackslash}p{0.25\textwidth}|>{\raggedright\arraybackslash}p{0.45\textwidth}|}
\hline
\textbf{Extending Use Case} & \textbf{Base Use Case} & \textbf{Condition} & \textbf{Added Behavior} \\ \hline
\endhead
E-01 Resend OTP & I-02 Verify Email OTP & The user requests a new OTP before verification. & Invalidate the previous OTP, generate a new code, and restart the 5-minute expiry period. \\ \hline
\end{longtable}
\endgroup
\par\medskip

\hypertarget{sec:a4-model-notes}{}
\addcontentsline{toc}{subsubsection}{A4. Model Notes}
\subsubsection*{A4. Model Notes}
\par\medskip

\begin{enumerate}
\item UC-01 is the selected general registration use case. G-01 and G-02 model the two registration paths without duplicating the base use case.
\item OTP verification is modeled under G-01 because it is the mandatory verification step for university-email registration. G-02 uses Google OAuth instead.
\item G-01 and G-02 inherit the UC-01 actor association and the I-01 validation behavior through use case generalization.
\end{enumerate}
\par\medskip

\hypertarget{sec:part-b-use-case-specification}{}
\addcontentsline{toc}{subsection}{Part B --- Use Case Specification}
\subsection*{Part B --- Use Case Specification}
\par\medskip

\hypertarget{sec:b1-brief-description}{}
\addcontentsline{toc}{subsubsection}{B1. Brief Description}
\subsubsection*{B1. Brief Description}
\par\medskip

A new university student creates a verified account and gains access to protected Studitopia features. The student can register with a supported university email address and a one-time password, or use Google OAuth. After successful registration and verification, the student can configure learning preferences and access the Home page.
\par\medskip

\hypertarget{sec:b2-preconditions}{}
\addcontentsline{toc}{subsubsection}{B2. Preconditions}
\subsubsection*{B2. Preconditions}
\par\medskip

\begin{enumerate}
\item The user does not have an active authenticated session.
\item The supported university and faculty email-domain catalog is available for university-email registration.
\item The email delivery service is available for the standard registration path.
\item Google OAuth 2.0 is configured for the optional Google registration path.
\end{enumerate}
\par\medskip

\hypertarget{sec:b3-postconditions}{}
\addcontentsline{toc}{subsubsection}{B3. Postconditions}
\subsubsection*{B3. Postconditions}
\par\medskip

\begin{enumerate}
\item \textbf{Success:} A user record exists with a unique username and email address, and the account is activated and email-verified.
\item \textbf{Success:} The user's password is stored only as a secure one-way hash; Google registration associates the authenticated Google identity with the account.
\item \textbf{Success:} The user is authorized to continue to Preference Settings and then the Home page.
\item \textbf{Failure:} No active account is created, and the user receives an error message that allows correction and retry.
\item \textbf{Failure:} An unverified account, if created, cannot access protected application features.
\end{enumerate}
\par\medskip

\hypertarget{sec:b4-main-success-scenario-basic-flow}{}
\addcontentsline{toc}{subsubsection}{B4. Main Success Scenario (Basic Flow)}
\subsubsection*{B4. Main Success Scenario (Basic Flow)}
\par\medskip

\begin{enumerate}
\item The user opens the registration page and selects \textbf{Register}.
\item The system displays Display Name, Username, Email, Password, Confirm Password, CAPTCHA, and a \textbf{Continue with Google} option.
\item The user enters the required details and completes the CAPTCHA.
\item The frontend validates that all required fields are present, the email format is valid, the username contains only alphanumeric characters, and the password satisfies the password policy.
\item The frontend verifies that Confirm Password matches Password.
\item The user submits the registration form.
\item The backend validates the submitted data, checks that the email and username are unique, and verifies that the email domain belongs to a supported university or faculty when university-email registration is used.
\item The backend stores a pending account with a securely hashed password and generates a time-limited OTP.
\item The email delivery service sends the OTP to the registered email address, and the system redirects the user to the OTP verification page.
\item The user enters the 6-digit OTP and submits it.
\item The backend verifies that the OTP is valid, unexpired, unused, and correctly formatted.
\item The system activates the account, updates the email-verification status, and authorizes access to Preference Settings.
\item The user completes or updates learning preferences and continues to the Home page.
\end{enumerate}
\par\medskip

\hypertarget{sec:b5-alternative-flows}{}
\addcontentsline{toc}{subsubsection}{B5. Alternative Flows}
\subsubsection*{B5. Alternative Flows}
\par\medskip

Unless otherwise stated, an alternative flow resumes at the next step of the basic flow.
\par\medskip

\hypertarget{sec:a1-google-registration}{}
\paragraph*{A1. Google registration}
\par\medskip

At basic flow step 2, the user selects \textbf{Continue with Google}. The system initiates Google OAuth, validates the authentication response, and creates or associates an account when registration requirements are satisfied. The user then continues to the appropriate next screen, normally Preference Settings.
\par\medskip

\hypertarget{sec:a2-invalid-client-side-input}{}
\paragraph*{A2. Invalid client-side input}
\par\medskip

At basic flow step 4 or 5, validation fails. The system displays field-level errors, keeps the entered values where appropriate, and does not submit the request until the errors are corrected.
\par\medskip

\hypertarget{sec:a3-duplicate-email-or-username}{}
\paragraph*{A3. Duplicate email or username}
\par\medskip

At basic flow step 7, the backend reports that the email or username is already registered. The system displays a clear error and offers navigation to Login or Password Recovery.
\par\medskip

\hypertarget{sec:a4-captcha-verification-failure}{}
\paragraph*{A4. CAPTCHA verification failure}
\par\medskip

At basic flow step 7, CAPTCHA verification fails. The system rejects the request, shows an error, and allows the user to retry.
\par\medskip

\hypertarget{sec:a5-resend-otp}{}
\paragraph*{A5. Resend OTP}
\par\medskip

At basic flow step 10, the user requests a new OTP. The system generates a new 6-digit OTP, invalidates the previous OTP, sends the new code, and restarts the 5-minute expiry period.
\par\medskip

\hypertarget{sec:a6-user-cancels-registration}{}
\paragraph*{A6. User cancels registration}
\par\medskip

The user leaves the registration flow before activation. No active account is available, and the pending verification data remains subject to the system's expiry and cleanup policy.
\par\medskip

\hypertarget{sec:b6-exception-flows}{}
\addcontentsline{toc}{subsubsection}{B6. Exception Flows}
\subsubsection*{B6. Exception Flows}
\par\medskip

\hypertarget{sec:e1-otp-invalid-expired-or-already-used}{}
\paragraph*{E1. OTP invalid, expired, or already used}
\par\medskip

The system rejects the OTP, displays an appropriate message, and keeps the account unverified. The user may request a new OTP.
\par\medskip

\hypertarget{sec:e2-otp-email-delivery-failure}{}
\paragraph*{E2. OTP email delivery failure}
\par\medskip

The account remains unverified and protected features remain unavailable. The system informs the user and allows another OTP request when permitted.
\par\medskip

\hypertarget{sec:e3-google-authentication-fails-or-is-cancelled}{}
\paragraph*{E3. Google authentication fails or is cancelled}
\par\medskip

No account is created or changed by the failed attempt. The system displays an error and allows the user to retry or use manual registration.
\par\medskip

\hypertarget{sec:e4-database-failure-during-account-creation}{}
\paragraph*{E4. Database failure during account creation}
\par\medskip

The registration transaction is rolled back, a generic error is shown, and the user can retry without creating a partial or duplicate account.
\par\medskip

\hypertarget{sec:e5-unverified-account-attempts-protected-access}{}
\paragraph*{E5. Unverified account attempts protected access}
\par\medskip

The backend denies access to protected features until email verification succeeds.
\par\medskip

\hypertarget{sec:b7-business-rules}{}
\addcontentsline{toc}{subsubsection}{B7. Business Rules}
\subsubsection*{B7. Business Rules}
\par\medskip

\begin{enumerate}
\item Display Name, Username, Email, Password, and Confirm Password are required.
\item A username must be unique and contain only alphanumeric characters.
\item An email address must be unique and valid; university-email registration requires a supported university or faculty domain.
\item A password must contain at least 8 characters, including at least one lowercase letter, one uppercase letter, one number, and one special character.
\item Passwords and Confirm Password fields are hidden by default and support a visibility toggle.
\item The OTP is a 6-digit numeric code, expires after 5 minutes, is single-use, and is invalidated when a new OTP is generated.
\item An account cannot access protected application features before successful email verification.
\item The password is persisted only as a secure one-way hash.
\end{enumerate}
\par\medskip

\hypertarget{sec:b8-special-requirements}{}
\addcontentsline{toc}{subsubsection}{B8. Special Requirements}
\subsubsection*{B8. Special Requirements}
\par\medskip

\begingroup
\small
\renewcommand{\baselinestretch}{1.0}\selectfont
\renewcommand{\arraystretch}{1.15}
\setlength{\tabcolsep}{6pt}
\begin{longtable}{|>{\raggedright\arraybackslash}p{0.10\textwidth}|>{\raggedright\arraybackslash}p{0.83\textwidth}|}
\hline
\textbf{ID} & \textbf{Requirement} \\ \hline
\endhead
NFR-SEC-2 & User credentials must be stored using one-way password hashing with a secure hashing algorithm. \\ \hline
NFR-SEC-3 & Sensitive communication must be encrypted using TLS 1.2 or higher. \\ \hline
NFR-SEC-8 & The system must protect against common web security vulnerabilities, including SQL injection, XSS, and CSRF. \\ \hline
NFR-USE-1 & A new university student must be able to complete registration within 5 minutes without external assistance. \\ \hline
Design constraint & The OTP is 6 digits, single-use, and expires after 5 minutes. The password policy and field visibility rules must be enforced by both frontend and backend validation. \\ \hline
\end{longtable}
\endgroup
\par\medskip

\hypertarget{sec:b9-extension-points}{}
\addcontentsline{toc}{subsubsection}{B9. Extension Points}
\subsubsection*{B9. Extension Points}
\par\medskip

\begingroup
\small
\renewcommand{\baselinestretch}{1.0}\selectfont
\renewcommand{\arraystretch}{1.15}
\setlength{\tabcolsep}{6pt}
\begin{longtable}{|>{\raggedright\arraybackslash}p{0.32\textwidth}|>{\raggedright\arraybackslash}p{0.61\textwidth}|}
\hline
\textbf{Extension Point} & \textbf{Location in Flow of Events} \\ \hline
\endhead
Google registration & Basic flow step 2; alternative flow A1 \\ \hline
Client-side validation recovery & Basic flow steps 4-5; alternative flow A2 \\ \hline
Duplicate-account handling & Basic flow step 7; alternative flow A3 \\ \hline
CAPTCHA retry & Basic flow step 7; alternative flow A4 \\ \hline
OTP resend & Basic flow step 10; alternative flow A5 and exceptions E1-E2 \\ \hline
Registration cancellation & Any step before activation; alternative flow A6 \\ \hline
Unverified-access recovery & Protected-feature access; exception E5 \\ \hline
\end{longtable}
\endgroup
\par\medskip

\hypertarget{sec:b10-interface-and-data-notes}{}
\addcontentsline{toc}{subsubsection}{B10. Interface and Data Notes}
\subsubsection*{B10. Interface and Data Notes}
\par\medskip

\begin{itemize}
\item \textbf{Screen/flow path:} \texttt{Authentication → Register → OTP Verification → Preference Settings → Home} (\texttt{feature\_flow.md}); \texttt{IsLogin = False → Authentication}, with \texttt{Login ↔ Register} and \texttt{Register → OTP → Preference Settings} (\texttt{screen\_flow.md}).
\item \textbf{Registration UI:} Display Name, Username, Email, Password, Confirm Password, CAPTCHA, \textbf{Register} submission, and \textbf{Continue with Google} option; Password and Confirm Password are hidden by default with a visibility toggle.
\item \textbf{OTP UI:} 6-slot numeric OTP input, masked recipient-email notice, \textbf{Resend code} control, expiry hint, and \textbf{Submit Button}; OTP feedback is shown on invalid or expired codes.
\item \textbf{Preference step:} During onboarding the user reaches Preference Settings after OTP verification and continues to Home; \texttt{settings.md} recommends \texttt{/settings/preferences} as the route.
\item \textbf{Documented API:} \texttt{POST /auth/register} with body \texttt{displayname}, \texttt{username}, \texttt{email}, \texttt{password}; \texttt{GET /auth/validate} is marked "Temporarily Disabled" and accepts optional \texttt{username}, \texttt{email}, \texttt{displayname}, \texttt{password} (not recommended) query parameters.
\item \textbf{Register responses:} \texttt{201 SUCCESS} with \texttt{access\_token} (\texttt{token}, \texttt{created}, \texttt{expired}); \texttt{400 BAD\_REQUEST}; \texttt{409 CONFLICT} with \texttt{DUPPLICATED\_USERNAME}, \texttt{DUPLICATED\_EMAIL}, \texttt{UNREGISTED\_EMAIL\_DOMAIN}; \texttt{422 UNPROCESSABLE} with \texttt{INVALID\_DISPLAYNAME}, \texttt{INVALID\_USERNAME}, \texttt{INVALID\_EMAIL}, \texttt{INVALID\_PASSWORD}; \texttt{500 INTERNAL\_SERVER\_ERROR} with \texttt{DATABASE\_ERROR}, \texttt{UNKNOWN\_ERROR}.
\item \textbf{Error codes:} \texttt{0 SUCCESS}, \texttt{2002 UNKNOWN ERROR}, \texttt{45000 UNKNOWN DATABASE ERROR}, \texttt{45001 INVALID\_DISPLAYNAME}, \texttt{45002 INVALID\_EMAIL}, \texttt{45003 INVALID\_USERNAME}, \texttt{45004 INVALID\_PASSWORD}, \texttt{46001 UNREGISTED\_EMAIL\_DOMAIN}, \texttt{47001 DUPLICATED\_EMAIL}, \texttt{47002 DUPLICATED\_USERNAME}.
\item \textbf{Account/status data:} \texttt{studitopia.register(displayname, username, email, password)} returns \texttt{session\_id}, \texttt{user\_id}, \texttt{status}; status values are \texttt{0 (inactive)}, \texttt{1 (preprofile)}, \texttt{2 (active)}, \texttt{3 (suspend)}. Account registered must be inactive; password must be hashed; HMAC-SHA256 is specified for \texttt{refresh\_token} and \texttt{access\_token}.
\item \textbf{Validation constraints:} Required fields; unique alphanumeric username; valid and unique email; supported university or faculty domain; minimum 8-character password with lowercase, uppercase, number, and special character; Confirm Password match; CAPTCHA verification on frontend and backend.
\item \textbf{OTP constraints:} 6-digit numeric code, sent to the registered email, stored as OTP or secure representation with expiration and verification status, expires after 5 minutes, single-use, invalidated on resend.
\item \textbf{Notifications and privacy:} OTP delivery is performed by the Email Delivery Service; sensitive communication uses TLS 1.2 or higher; password hashes are never displayed; unverified accounts are restricted from protected application features.
\end{itemize}
\par\medskip

\hypertarget{sec:b11-source-alignment-notes}{}
\addcontentsline{toc}{subsubsection}{B11. Source Alignment Notes}
\subsubsection*{B11. Source Alignment Notes}
\par\medskip

\begin{itemize}
\item API coverage is partial: no endpoint is documented for OTP submission, OTP resend, Google OAuth registration, or the Preference Settings transition; only \texttt{POST /auth/register} and the temporarily disabled \texttt{GET /auth/validate} are present.
\item \texttt{POST /auth/register} returns an \texttt{access\_token}, while the API also requires the registered account to be inactive and UC-01 postconditions forbid protected access before email verification; the intended token scope is not specified.
\item CAPTCHA verification is required by \texttt{FR-01.17} and business rules, but no dedicated CAPTCHA error code is listed in \texttt{response\_code.md}.
\item Error-name spelling differs across sources: \texttt{API.md} uses \texttt{DUPPLICATED\_EMAIL}, while \texttt{response\_code.md} uses \texttt{DUPLICATED\_EMAIL}; the unusual \texttt{UNREGISTED\_EMAIL\_DOMAIN} spelling is preserved from both.
\item The focused model states G-02 uses Google OAuth instead of OTP, but the sources do not explicitly state whether the Google-authenticated email is treated as verified; A1 only says the user continues, normally to Preference Settings.
\item \texttt{/settings/preferences} is a recommended route in \texttt{settings.md}, not a route defined in the proposal flow documents.
\end{itemize}
\par\medskip

\par\medskip
\hypertarget{sec:ui-prototype-specimens}{}
\addcontentsline{toc}{subsection}{UI Prototype Specimens}
\subsection*{UI Prototype Specimens}
\par\medskip
\begin{figure}[H]
\centering
\includegraphics[width=0.75\textwidth,height=0.42\textheight,keepaspectratio]{"../../prototype/otp_desktop.png"}
\par\vspace{2pt}
{\footnotesize \textbf{UI Prototype:} OTP - Desktop}
\end{figure}
\par\smallskip
\begin{figure}[H]
\centering
\includegraphics[width=0.75\textwidth,height=0.42\textheight,keepaspectratio]{"../../prototype/register_desktop.png"}
\par\vspace{2pt}
{\footnotesize \textbf{UI Prototype:} Register - Desktop}
\end{figure}
\par\smallskip

\pagebreak
\hypertarget{sec:uc-02-log-in}{}
\addcontentsline{toc}{section}{UC-02 --- Log In}
\section*{UC-02 --- Log In}
\par\medskip

\textbf{Version:} \texttt{2.0}
\par\medskip

\textbf{Performed by:} [Nguyen Le Quoc Huy] | \textbf{Edited by:} [Nguyen Thanh An]
\par\medskip

\begingroup
\small
\renewcommand{\baselinestretch}{1.0}\selectfont
\renewcommand{\arraystretch}{1.15}
\setlength{\tabcolsep}{6pt}
\begin{longtable}{|>{\raggedright\arraybackslash}p{0.09\textwidth}|>{\raggedright\arraybackslash}p{0.84\textwidth}|}
\hline
\textbf{Field} & \textbf{Value} \\ \hline
\endhead
Use Case ID & UC-02 \\ \hline
Name & Log In \\ \hline
Primary Actor & Registered University Student \\ \hline
Supporting Actors & Google OAuth 2.0 \\ \hline
Feature Group & Authentication \\ \hline
Priority & High \\ \hline
Related User Stories & US-02 \\ \hline
Related Requirements & US-02, US-04; FR-02.1 to FR-02.5, FR-04.1 to FR-04.4; NFR-SEC-1, NFR-SEC-3, NFR-SEC-6, NFR-SEC-7, NFR-SEC-8 \\ \hline
Source Documents & \texttt{use\_cases\_model/UC-02-log-in.md}; \texttt{use\_cases\_description.md} (UC-02); \texttt{use\_cases\_model.md}; \texttt{use\_case\_diagrams.md}; \texttt{PA3-Group05/use\_case.md} (§2 UC-02); \texttt{overview.md}; \texttt{functional\_requirement.md}; \texttt{non\_functional\_requirement.md}; \texttt{userstory.md}; \texttt{feature\_flow.md}; \texttt{screen\_flow.md}; \texttt{API.md}; \texttt{response\_code.md}; \texttt{settings.md} \\ \hline
\end{longtable}
\endgroup
\par\medskip

\hypertarget{sec:part-a-use-case-model}{}
\addcontentsline{toc}{subsection}{Part A --- Use Case Model}
\subsection*{Part A --- Use Case Model}
\par\medskip

\hypertarget{sec:a1-focused-use-case-diagram}{}
\addcontentsline{toc}{subsubsection}{A1. Focused Use Case Diagram}
\subsubsection*{A1. Focused Use Case Diagram}
\par\medskip

\begin{center}
\includegraphics[width=0.85\textwidth,height=0.75\textheight,keepaspectratio]{"figures/UC-02-log-in_1.png"}
\end{center}
\par\medskip

\textbf{Diagram notation.} Person shape = human actor; cloud = supporting system; blue = selected base use case; purple = included; green = extending; solid line = association; dashed arrow = include or extend; labelled solid arrow = generalization.
\par\medskip

\hypertarget{sec:a2-actor-register}{}
\addcontentsline{toc}{subsubsection}{A2. Actor Register}
\subsubsection*{A2. Actor Register}
\par\medskip

\begingroup
\small
\renewcommand{\baselinestretch}{1.0}\selectfont
\renewcommand{\arraystretch}{1.15}
\setlength{\tabcolsep}{6pt}
\begin{longtable}{|>{\raggedright\arraybackslash}p{0.16\textwidth}|>{\raggedright\arraybackslash}p{0.10\textwidth}|>{\raggedright\arraybackslash}p{0.66\textwidth}|}
\hline
\textbf{Actor} & \textbf{Type} & \textbf{Description} \\ \hline
\endhead
Registered University Student & Primary & A student with an existing account who logs in to access protected features. \\ \hline
University Student & General human actor & Parent of Registered University Student; inherits general student associations to UC-03, UC-04, UC-05, UC-06, and UC-07. \\ \hline
Google OAuth 2.0 & Supporting System & Provides an optional external authentication and registration path; supports E-02 Log In with Google. \\ \hline
\end{longtable}
\endgroup
\par\medskip

\hypertarget{sec:a3-relationships}{}
\addcontentsline{toc}{subsubsection}{A3. Relationships}
\subsubsection*{A3. Relationships}
\par\medskip

\hypertarget{sec:a3-1-associations}{}
\paragraph*{A3.1 Associations}
\par\medskip

\begingroup
\small
\renewcommand{\baselinestretch}{1.0}\selectfont
\renewcommand{\arraystretch}{1.15}
\setlength{\tabcolsep}{6pt}
\begin{longtable}{|>{\raggedright\arraybackslash}p{0.21\textwidth}|>{\raggedright\arraybackslash}p{0.17\textwidth}|>{\raggedright\arraybackslash}p{0.10\textwidth}|>{\raggedright\arraybackslash}p{0.44\textwidth}|}
\hline
\textbf{Actor} & \textbf{Use Case} & \textbf{Type} & \textbf{Description} \\ \hline
\endhead
Registered University Student & UC-02 Log In & Primary & Authenticates with credentials and reaches protected features. \\ \hline
Google OAuth 2.0 & E-02 Log In with Google & Supporting & Authenticates the optional Google login path. \\ \hline
\end{longtable}
\endgroup
\par\medskip

\hypertarget{sec:a3-2-generalization-and-specialization}{}
\paragraph*{A3.2 Generalization and Specialization}
\par\medskip

\begingroup
\small
\renewcommand{\baselinestretch}{1.0}\selectfont
\renewcommand{\arraystretch}{1.15}
\setlength{\tabcolsep}{6pt}
\begin{longtable}{|>{\raggedright\arraybackslash}p{0.21\textwidth}|>{\raggedright\arraybackslash}p{0.13\textwidth}|>{\raggedright\arraybackslash}p{0.15\textwidth}|>{\raggedright\arraybackslash}p{0.42\textwidth}|}
\hline
\textbf{Child} & \textbf{Parent} & \textbf{Type} & \textbf{Inherited Behavior or Association} \\ \hline
\endhead
Registered University Student & University Student & Actor generalization & Inherits the general student associations and adds UC-02. \\ \hline
\end{longtable}
\endgroup
\par\medskip

\hypertarget{sec:a3-3-include}{}
\paragraph*{A3.3 Include}
\par\medskip

\begingroup
\small
\renewcommand{\baselinestretch}{1.0}\selectfont
\renewcommand{\arraystretch}{1.15}
\setlength{\tabcolsep}{6pt}
\begin{longtable}{|>{\raggedright\arraybackslash}p{0.10\textwidth}|>{\raggedright\arraybackslash}p{0.26\textwidth}|>{\raggedright\arraybackslash}p{0.56\textwidth}|}
\hline
\textbf{Base Use Case} & \textbf{Included Use Case} & \textbf{Meaning} \\ \hline
\endhead
UC-02 Log In & I-03 Create Authenticated Session & Create the authenticated session and return the student's access context. \\ \hline
\end{longtable}
\endgroup
\par\medskip

\hypertarget{sec:a3-4-extend}{}
\paragraph*{A3.4 Extend}
\par\medskip

\begingroup
\small
\renewcommand{\baselinestretch}{1.0}\selectfont
\renewcommand{\arraystretch}{1.15}
\setlength{\tabcolsep}{6pt}
\begin{longtable}{|>{\raggedright\arraybackslash}p{0.16\textwidth}|>{\raggedright\arraybackslash}p{0.10\textwidth}|>{\raggedright\arraybackslash}p{0.38\textwidth}|>{\raggedright\arraybackslash}p{0.28\textwidth}|}
\hline
\textbf{Extending Use Case} & \textbf{Base Use Case} & \textbf{Condition} & \textbf{Added Behavior} \\ \hline
\endhead
E-02 Log In with Google & UC-02 Log In & The student selects Google login instead of credentials. & Add the Google OAuth authentication path. \\ \hline
\end{longtable}
\endgroup
\par\medskip

\hypertarget{sec:a4-model-notes}{}
\addcontentsline{toc}{subsubsection}{A4. Model Notes}
\subsubsection*{A4. Model Notes}
\par\medskip

\begin{enumerate}
\item Successful login always includes I-03 because a valid authenticated session is required before protected features are available.
\item Google login is an optional extension of UC-02. The standard username-or-email and password path remains the base behavior.
\item Registered University Student inherits all general University Student associations in addition to the direct UC-02 association.
\end{enumerate}
\par\medskip

\hypertarget{sec:part-b-use-case-specification}{}
\addcontentsline{toc}{subsection}{Part B --- Use Case Specification}
\subsection*{Part B --- Use Case Specification}
\par\medskip

\hypertarget{sec:b1-brief-description}{}
\addcontentsline{toc}{subsubsection}{B1. Brief Description}
\subsubsection*{B1. Brief Description}
\par\medskip

A registered university student authenticates and regains access to their profile, university forum, posts, and learning activities. Login supports a registered username or university email with a password, and an optional Google OAuth path.
\par\medskip

\hypertarget{sec:b2-preconditions}{}
\addcontentsline{toc}{subsubsection}{B2. Preconditions}
\subsubsection*{B2. Preconditions}
\par\medskip

\begin{enumerate}
\item The student has a registered and activated account.
\item The student is not already authenticated, or chooses to start a new authenticated session.
\item The authentication service and, when used, Google OAuth 2.0 are available.
\end{enumerate}
\par\medskip

\hypertarget{sec:b3-postconditions}{}
\addcontentsline{toc}{subsubsection}{B3. Postconditions}
\subsubsection*{B3. Postconditions}
\par\medskip

\begin{enumerate}
\item \textbf{Success:} A valid authenticated session is created and the student reaches the Home page or the originally requested protected page.
\item \textbf{Success:} The student's default university forum context is available after login.
\item \textbf{Failure:} No authenticated session is created, and the student remains on the login page with a non-disclosing error message.
\end{enumerate}
\par\medskip

\hypertarget{sec:b4-main-success-scenario-basic-flow}{}
\addcontentsline{toc}{subsubsection}{B4. Main Success Scenario (Basic Flow)}
\subsubsection*{B4. Main Success Scenario (Basic Flow)}
\par\medskip

\begin{enumerate}
\item The student opens the login page.
\item The system displays the username-or-email field, password field, CAPTCHA, \textbf{Remember Me} option where supported, \textbf{Forget Password}, \textbf{Register}, and \textbf{Continue with Google} options.
\item The student enters a registered username or university email address and password, then completes the CAPTCHA.
\item The frontend performs basic input validation and submits the login request.
\item The backend verifies the CAPTCHA, credentials, account activation status, and any account restrictions.
\item The backend creates an authenticated session and returns the access result.
\item The frontend stores the session according to the \textbf{Remember Me} choice, loads the student's profile and forum context, and redirects to Home or the intended protected page.
\item The system displays the Home page with the Forum entry and the student's university forum as the default context.
\end{enumerate}
\par\medskip

\hypertarget{sec:b5-alternative-flows}{}
\addcontentsline{toc}{subsubsection}{B5. Alternative Flows}
\subsubsection*{B5. Alternative Flows}
\par\medskip

Unless otherwise stated, an alternative flow resumes at the next step of the basic flow.
\par\medskip

\hypertarget{sec:a1-google-login}{}
\paragraph*{A1. Google login}
\par\medskip

At basic flow step 2, the student selects \textbf{Continue with Google}. The system completes Google OAuth, validates the response, and either creates the authenticated session or displays an appropriate error.
\par\medskip

\hypertarget{sec:a2-remember-me-selected}{}
\paragraph*{A2. Remember Me selected}
\par\medskip

At basic flow step 7, the system maintains the session across browser restarts according to the selected option and the platform session policy.
\par\medskip

\hypertarget{sec:a3-intended-page-available}{}
\paragraph*{A3. Intended page available}
\par\medskip

If the student was redirected from a protected page, the system returns the student to that page after successful authentication instead of Home.
\par\medskip

\hypertarget{sec:a4-unverified-account}{}
\paragraph*{A4. Unverified account}
\par\medskip

If the backend finds that the account is not verified, login is denied and the student is directed to complete OTP verification.
\par\medskip

\hypertarget{sec:b6-exception-flows}{}
\addcontentsline{toc}{subsubsection}{B6. Exception Flows}
\subsubsection*{B6. Exception Flows}
\par\medskip

\hypertarget{sec:e1-invalid-credentials}{}
\paragraph*{E1. Invalid credentials}
\par\medskip

The system displays a clear error without revealing whether the username, email, or password was incorrect.
\par\medskip

\hypertarget{sec:e2-captcha-failure}{}
\paragraph*{E2. CAPTCHA failure}
\par\medskip

The system rejects the request, displays an error, and allows the student to retry.
\par\medskip

\hypertarget{sec:e3-unavailable-or-restricted-account}{}
\paragraph*{E3. Unavailable or restricted account}
\par\medskip

The system denies authentication and displays an appropriate message without exposing sensitive account details.
\par\medskip

\hypertarget{sec:e4-authentication-service-failure}{}
\paragraph*{E4. Authentication service failure}
\par\medskip

No session is created. The system displays a temporary error and allows the student to retry.
\par\medskip

\hypertarget{sec:e5-session-expires-while-accessing-a-protected-page}{}
\paragraph*{E5. Session expires while accessing a protected page}
\par\medskip

The system denies access, informs the student that the session expired, and requires login again.
\par\medskip

\hypertarget{sec:b7-business-rules}{}
\addcontentsline{toc}{subsubsection}{B7. Business Rules}
\subsubsection*{B7. Business Rules}
\par\medskip

\begin{enumerate}
\item Login is accepted with either a registered username or a registered university email address.
\item Invalid-credential messages must not reveal which credential field was incorrect.
\item CAPTCHA must be successfully verified before login processing.
\item Google login is an optional authentication path for registered users.
\item Protected pages and APIs require a valid authenticated session.
\item After successful login, the student's profile, forum, posts, and learning activities are accessible according to their permissions.
\item The session is subject to the platform's inactivity-expiry policy.
\end{enumerate}
\par\medskip

\hypertarget{sec:b8-special-requirements}{}
\addcontentsline{toc}{subsubsection}{B8. Special Requirements}
\subsubsection*{B8. Special Requirements}
\par\medskip

\begingroup
\small
\renewcommand{\baselinestretch}{1.0}\selectfont
\renewcommand{\arraystretch}{1.15}
\setlength{\tabcolsep}{6pt}
\begin{longtable}{|>{\raggedright\arraybackslash}p{0.10\textwidth}|>{\raggedright\arraybackslash}p{0.83\textwidth}|}
\hline
\textbf{ID} & \textbf{Requirement} \\ \hline
\endhead
NFR-SEC-1 & Authentication is required before accessing protected user resources. \\ \hline
NFR-SEC-3 & Sensitive communication must be encrypted using TLS 1.2 or higher. \\ \hline
NFR-SEC-6 & Server-side authorization checks must prevent unauthorized access. \\ \hline
NFR-SEC-7 & Inactive authentication sessions must expire automatically after 30 minutes. \\ \hline
NFR-SEC-8 & The system must protect against SQL injection, XSS, CSRF, and other common web vulnerabilities. \\ \hline
\end{longtable}
\endgroup
\par\medskip

\hypertarget{sec:b9-extension-points}{}
\addcontentsline{toc}{subsubsection}{B9. Extension Points}
\subsubsection*{B9. Extension Points}
\par\medskip

\begingroup
\small
\renewcommand{\baselinestretch}{1.0}\selectfont
\renewcommand{\arraystretch}{1.15}
\setlength{\tabcolsep}{6pt}
\begin{longtable}{|>{\raggedright\arraybackslash}p{0.39\textwidth}|>{\raggedright\arraybackslash}p{0.54\textwidth}|}
\hline
\textbf{Extension Point} & \textbf{Location in Flow of Events} \\ \hline
\endhead
Google login & Basic flow step 2; alternative flow A1 \\ \hline
Remember Me & Basic flow step 7; alternative flow A2 \\ \hline
Return to intended page & Basic flow step 8; alternative flow A3 \\ \hline
Unverified-account recovery & Basic flow step 5; alternative flow A4 \\ \hline
Session expiry & Protected-page access; exception E5 \\ \hline
\end{longtable}
\endgroup
\par\medskip

\hypertarget{sec:b10-interface-and-data-notes}{}
\addcontentsline{toc}{subsubsection}{B10. Interface and Data Notes}
\subsubsection*{B10. Interface and Data Notes}
\par\medskip

\begin{itemize}
\item Screen/flow path: \texttt{Start → IsLogin (False) → Authentication → Login}; successful \texttt{Login → Home}; \texttt{Login → Register}, \texttt{Login → Recover}, \texttt{Login → Home}; \texttt{Home → Settings → Logout → Login}; feature flow \texttt{Authentication → Login → Home}.
\item Login UI fields and actions: Username / Email, Password, CAPTCHA, \textbf{Remember Me} where supported, \textbf{Forget Password}, \textbf{Register}, and \textbf{Continue with Google}.
\item \texttt{POST /auth/login} query \texttt{type: 'token' | 'username' | 'email'}; token body \texttt{\{\}}; username/email body \texttt{\{ identifier, password \}}.
\item \texttt{POST /auth/login} responses: \texttt{successCode: 200} for \texttt{token} with \texttt{access\_token}, \texttt{created}, \texttt{expired}; \texttt{successCode: 201} for \texttt{username} or \texttt{email} with \texttt{access\_token: \{ token, created, expired \}}; \texttt{BAD\_REQUEST} 400; \texttt{UNAUTHORIZED} 401; \texttt{CONFLICT} 409 (\texttt{DUPPLICATED\_USERNAME | DUPLICATED\_EMAIL | UNREGISTED\_EMAIL\_DOMAIN}); \texttt{UNPROCESSABLE} 422 (\texttt{INVALID\_USERNAME | INVALID\_EMAIL | INVALID\_PASSWORD}); \texttt{INTERNAL\_SERVER\_ERROR} 500 (\texttt{DATABASE\_ERROR | UNKNOWN\_ERROR}).
\item Error-code catalog: \texttt{UNKNOWN ERROR} 45000, \texttt{INVALID\_EMAIL} 45002, \texttt{INVALID\_USERNAME} 45003, \texttt{INVALID\_PASSWORD} 45004, \texttt{UNREGISTED\_EMAIL\_DOMAIN} 46001, \texttt{DUPLICATED\_EMAIL} 47001, \texttt{DUPLICATED\_USERNAME} 47002; backend \texttt{UNKNOWN ERROR} 2002 and \texttt{UNKNOWN DATABASE ERROR} 45000.
\item \texttt{GET /auth/validate} is listed as "Temporarily Disabled" with responses 204, 400, 409, 422, and 500; it is not the main login submission endpoint.
\item Session/Remember Me: \textbf{Remember Me} maintains the session across browser restarts (AC-02.3); expired sessions are denied on protected pages/APIs (FR-04.2), the inactivity timer resets on a qualifying request (FR-04.3), and expiration requires re-login (FR-04.4).
\item Data constraints: login accepts a registered username or university email address; CAPTCHA must be verified; backend checks activation status and restrictions; API account status values are inactive (0), preprofile (1), active (2), suspend (3).
\end{itemize}
\par\medskip

\hypertarget{sec:b11-source-alignment-notes}{}
\addcontentsline{toc}{subsubsection}{B11. Source Alignment Notes}
\subsubsection*{B11. Source Alignment Notes}
\par\medskip

\begin{itemize}
\item Inactivity timeout conflicts: FR-04.1 and US-04/AC-04.1 state 15 minutes, while NFR-SEC-7 states 30 minutes; \texttt{settings.md} recommends the stricter 15-minute value unless the product explicitly decides otherwise.
\item \texttt{API.md} defines \texttt{POST /auth/login} but no logout endpoint, although US-03 and FR-03.1 to FR-03.3 require logout, client cleanup, and backend session invalidation.
\item CAPTCHA is required by AC-02.4 and the use-case flow, but the documented \texttt{POST /auth/login} request body lists only \texttt{identifier} and \texttt{password} for username/email login.
\item Google login is required by FR-02.4 and AC-02.5 and modeled as E-02, but \texttt{API.md} does not define a dedicated Google OAuth login endpoint.
\end{itemize}
\par\medskip

\par\medskip
\hypertarget{sec:ui-prototype-specimens}{}
\addcontentsline{toc}{subsection}{UI Prototype Specimens}
\subsection*{UI Prototype Specimens}
\par\medskip
\begin{figure}[H]
\centering
\includegraphics[width=0.75\textwidth,height=0.42\textheight,keepaspectratio]{"../../prototype/login_desktop.png"}
\par\vspace{2pt}
{\footnotesize \textbf{UI Prototype:} Login - Desktop}
\end{figure}
\par\smallskip
\begin{figure}[H]
\centering
\includegraphics[width=0.75\textwidth,height=0.42\textheight,keepaspectratio]{"../../prototype/settings_login_desktop_destination.png"}
\par\vspace{2pt}
{\footnotesize \textbf{UI Prototype:} Settings: Login - Desktop destination}
\end{figure}
\par\smallskip

\pagebreak
\hypertarget{sec:uc-08-logout}{}
\addcontentsline{toc}{section}{UC-08 --- Logout}
\section*{UC-08 --- Logout}
\par\medskip

\textbf{Version:} \texttt{1.0}
\par\medskip

\textbf{Performed by:} [Nguyen Thanh An] | \textbf{Edited by:} [Nguyen Thanh An]
\par\medskip

\begingroup
\small
\renewcommand{\baselinestretch}{1.0}\selectfont
\renewcommand{\arraystretch}{1.15}
\setlength{\tabcolsep}{6pt}
\begin{longtable}{|>{\raggedright\arraybackslash}p{0.10\textwidth}|>{\raggedright\arraybackslash}p{0.83\textwidth}|}
\hline
\textbf{Field} & \textbf{Value} \\ \hline
\endhead
Use Case ID & UC-08 \\ \hline
Name & Logout \\ \hline
Primary Actor & logged-in user \\ \hline
Supporting Actors & None \\ \hline
Feature Group & Authentication Account \\ \hline
Priority & Not explicitly stated \\ \hline
Related User Stories & US-03 \\ \hline
Related Requirements & FR-03.1, FR-03.2, FR-03.3; AC-03.1, AC-03.2, AC-03.3, AC-03.4; NFR-SEC-1, NFR-SEC-3, NFR-SEC-6, NFR-SEC-8 \\ \hline
Source Documents & \texttt{proposals/userstory.md}; \texttt{proposals/functional\_requirement.md}; \texttt{proposals/overview.md}; \texttt{proposals/feature\_flow.md}; \texttt{proposals/screen\_flow.md}; \texttt{proposals/API.md}; \texttt{proposals/response\_code.md}; \texttt{proposals/non\_functional\_requirement.md}; \texttt{proposals/settings.md} \\ \hline
\end{longtable}
\endgroup
\par\medskip

\hypertarget{sec:part-a-use-case-model}{}
\addcontentsline{toc}{subsection}{Part A --- Use Case Model}
\subsection*{Part A --- Use Case Model}
\par\medskip

\hypertarget{sec:a1-focused-use-case-diagram}{}
\addcontentsline{toc}{subsubsection}{A1. Focused Use Case Diagram}
\subsubsection*{A1. Focused Use Case Diagram}
\par\medskip

\begin{center}
\includegraphics[width=0.85\textwidth,height=0.75\textheight,keepaspectratio]{"figures/UC-08-logout_1.png"}
\end{center}
\par\medskip

\textbf{Diagram notation.} Person shape = human actor; cloud = supporting system; blue = selected use case; solid line = association; labelled solid arrow = generalization; dashed arrow = include or extend.
\par\medskip

\hypertarget{sec:a2-actor-register}{}
\addcontentsline{toc}{subsubsection}{A2. Actor Register}
\subsubsection*{A2. Actor Register}
\par\medskip

\begingroup
\small
\renewcommand{\baselinestretch}{1.0}\selectfont
\renewcommand{\arraystretch}{1.15}
\setlength{\tabcolsep}{6pt}
\begin{longtable}{|>{\raggedright\arraybackslash}p{0.10\textwidth}|>{\raggedright\arraybackslash}p{0.10\textwidth}|>{\raggedright\arraybackslash}p{0.73\textwidth}|}
\hline
\textbf{Actor} & \textbf{Type} & \textbf{Description} \\ \hline
\endhead
logged-in user & Primary human actor & A user with an active authenticated session who ends that session to protect account security on a shared device or when finished using the platform. \\ \hline
\end{longtable}
\endgroup
\par\medskip

\hypertarget{sec:a3-relationships}{}
\addcontentsline{toc}{subsubsection}{A3. Relationships}
\subsubsection*{A3. Relationships}
\par\medskip

\hypertarget{sec:a3-1-associations}{}
\paragraph*{A3.1 Associations}
\par\medskip

\begingroup
\small
\renewcommand{\baselinestretch}{1.0}\selectfont
\renewcommand{\arraystretch}{1.15}
\setlength{\tabcolsep}{6pt}
\begin{longtable}{|>{\raggedright\arraybackslash}p{0.10\textwidth}|>{\raggedright\arraybackslash}p{0.09\textwidth}|>{\raggedright\arraybackslash}p{0.09\textwidth}|>{\raggedright\arraybackslash}p{0.62\textwidth}|}
\hline
\textbf{Actor} & \textbf{Use Case} & \textbf{Type} & \textbf{Description} \\ \hline
\endhead
logged-in user & UC-08 Logout & Primary & Selects the logout action and causes client cleanup and backend session invalidation. \\ \hline
\end{longtable}
\endgroup
\par\medskip

\hypertarget{sec:a3-2-generalization-and-specialization}{}
\paragraph*{A3.2 Generalization and Specialization}
\par\medskip

\begingroup
\small
\renewcommand{\baselinestretch}{1.0}\selectfont
\renewcommand{\arraystretch}{1.15}
\setlength{\tabcolsep}{6pt}
\begin{longtable}{|>{\raggedright\arraybackslash}p{0.09\textwidth}|>{\raggedright\arraybackslash}p{0.09\textwidth}|>{\raggedright\arraybackslash}p{0.09\textwidth}|>{\raggedright\arraybackslash}p{0.64\textwidth}|}
\hline
\textbf{Child} & \textbf{Parent} & \textbf{Type} & \textbf{Inherited Behavior or Association} \\ \hline
\endhead
None & None & None & No actor or use-case generalization is specified in the proposals. \\ \hline
\end{longtable}
\endgroup
\par\medskip

\hypertarget{sec:a3-3-include}{}
\paragraph*{A3.3 Include}
\par\medskip

\begingroup
\small
\renewcommand{\baselinestretch}{1.0}\selectfont
\renewcommand{\arraystretch}{1.15}
\setlength{\tabcolsep}{6pt}
\begin{longtable}{|>{\raggedright\arraybackslash}p{0.16\textwidth}|>{\raggedright\arraybackslash}p{0.20\textwidth}|>{\raggedright\arraybackslash}p{0.56\textwidth}|}
\hline
\textbf{Base Use Case} & \textbf{Included Use Case} & \textbf{Meaning} \\ \hline
\endhead
None & None & No included use case is named in the proposals. \\ \hline
\end{longtable}
\endgroup
\par\medskip

\hypertarget{sec:a3-4-extend}{}
\paragraph*{A3.4 Extend}
\par\medskip

\begingroup
\small
\renewcommand{\baselinestretch}{1.0}\selectfont
\renewcommand{\arraystretch}{1.15}
\setlength{\tabcolsep}{6pt}
\begin{longtable}{|>{\raggedright\arraybackslash}p{0.19\textwidth}|>{\raggedright\arraybackslash}p{0.13\textwidth}|>{\raggedright\arraybackslash}p{0.10\textwidth}|>{\raggedright\arraybackslash}p{0.49\textwidth}|}
\hline
\textbf{Extending Use Case} & \textbf{Base Use Case} & \textbf{Condition} & \textbf{Added Behavior} \\ \hline
\endhead
None & None & None & No extending use case is named in the proposals. \\ \hline
\end{longtable}
\endgroup
\par\medskip

\hypertarget{sec:a4-model-notes}{}
\addcontentsline{toc}{subsubsection}{A4. Model Notes}
\subsubsection*{A4. Model Notes}
\par\medskip

\begin{enumerate}
\item UC-08 is limited to ending the current authenticated session; it does not model login, registration, password recovery, or account deletion.
\item Client cleanup and backend invalidation are internal behavior of UC-08 because the proposals do not name a separate included use case.
\item The logout action is available after navigation to Settings or from global navigation, but no distinct UI actor or supporting service is specified.
\item No new I-xx, E-xx, or G-xx identifiers are introduced; the model contains only the primary actor and the selected base use case.
\end{enumerate}
\par\medskip

\hypertarget{sec:part-b-use-case-specification}{}
\addcontentsline{toc}{subsection}{Part B --- Use Case Specification}
\subsection*{Part B --- Use Case Specification}
\par\medskip

\hypertarget{sec:b1-brief-description}{}
\addcontentsline{toc}{subsubsection}{B1. Brief Description}
\subsubsection*{B1. Brief Description}
\par\medskip

A logged-in user ends the current session to protect the account when using a shared device or when finished with the platform. The system clears authentication tokens and temporary user data from local and session storage and invalidates the active backend session.
\par\medskip

The screen flow places logout after Home or Settings and returns the user to Login. Settings groups the control as a global account action labeled \textbf{Log out}.
\par\medskip

\hypertarget{sec:b2-preconditions}{}
\addcontentsline{toc}{subsubsection}{B2. Preconditions}
\subsubsection*{B2. Preconditions}
\par\medskip

\begin{enumerate}
\item The user has an active authenticated session.
\item The user can reach the logout control from Settings or global navigation.
\item The frontend and backend are available to perform cleanup and session invalidation.
\end{enumerate}
\par\medskip

\hypertarget{sec:b3-postconditions}{}
\addcontentsline{toc}{subsubsection}{B3. Postconditions}
\subsubsection*{B3. Postconditions}
\par\medskip

\begin{enumerate}
\item \textbf{Success:} Authentication tokens and temporary user data are cleared from local and session storage.
\item \textbf{Success:} The active backend session is invalidated and the user is returned to Login.
\item \textbf{Failure:} If cleanup or invalidation does not complete, the source documents do not specify a recovery message or retries.
\end{enumerate}
\par\medskip

\hypertarget{sec:b4-main-success-scenario-basic-flow}{}
\addcontentsline{toc}{subsubsection}{B4. Main Success Scenario (Basic Flow)}
\subsubsection*{B4. Main Success Scenario (Basic Flow)}
\par\medskip

\begin{enumerate}
\item The logged-in user opens Settings or uses global navigation.
\item The system displays the \textbf{Log out} control.
\item The user selects \textbf{Log out}.
\item The frontend clears authentication tokens and temporary user data from local and session storage.
\item The backend invalidates the user's active session.
\item The system redirects the user to Login.
\end{enumerate}
\par\medskip

\hypertarget{sec:b5-alternative-flows}{}
\addcontentsline{toc}{subsubsection}{B5. Alternative Flows}
\subsubsection*{B5. Alternative Flows}
\par\medskip

The story specifies no distinct alternative flow; the source-backed variants are the entry points to the same logout behavior.
\par\medskip

\hypertarget{sec:a1-logout-from-settings}{}
\paragraph*{A1. Logout from Settings}
\par\medskip

From \texttt{Home → Settings}, the user selects \textbf{Log out} in the Settings area. The flow resumes at basic flow step 3.
\par\medskip

\hypertarget{sec:a2-logout-from-global-navigation}{}
\paragraph*{A2. Logout from global navigation}
\par\medskip

The user selects \textbf{Log out} from global navigation without opening Settings. The flow resumes at basic flow step 3.
\par\medskip

\hypertarget{sec:b6-exception-flows}{}
\addcontentsline{toc}{subsubsection}{B6. Exception Flows}
\subsubsection*{B6. Exception Flows}
\par\medskip

No logout-specific exception flow is specified in the proposals.
\par\medskip

\hypertarget{sec:b7-business-rules}{}
\addcontentsline{toc}{subsubsection}{B7. Business Rules}
\subsubsection*{B7. Business Rules}
\par\medskip

\begin{enumerate}
\item A logged-in user is allowed to log out of the account (FR-03.1).
\item On logout, all authentication tokens and temporary user data stored in local and session storage must be cleared (FR-03.2, AC-03.1).
\item On logout, the user's active session must be invalidated on the backend (FR-03.3, AC-03.2).
\item UI must appear exactly as designed, and navigation buttons must work correctly (AC-03.3, AC-03.4).
\item After logout, protected user resources remain subject to authentication requirements (NFR-SEC-1).
\end{enumerate}
\par\medskip

\hypertarget{sec:b8-special-requirements}{}
\addcontentsline{toc}{subsubsection}{B8. Special Requirements}
\subsubsection*{B8. Special Requirements}
\par\medskip

\begingroup
\small
\renewcommand{\baselinestretch}{1.0}\selectfont
\renewcommand{\arraystretch}{1.15}
\setlength{\tabcolsep}{6pt}
\begin{longtable}{|>{\raggedright\arraybackslash}p{0.11\textwidth}|>{\raggedright\arraybackslash}p{0.82\textwidth}|}
\hline
\textbf{ID} & \textbf{Requirement} \\ \hline
\endhead
NFR-SEC-1 & Authentication is required before accessing protected user resources. \\ \hline
NFR-SEC-3 & Sensitive communication must be encrypted using TLS 1.2 or higher. \\ \hline
NFR-SEC-6 & Server-side authorization checks must prevent unauthorized access. \\ \hline
NFR-SEC-8 & The system must protect against common web security vulnerabilities, including SQL injection, XSS, and CSRF. \\ \hline
Design constraint & Local and session storage cleanup and backend session invalidation are required by FR-03.2, FR-03.3, AC-03.1, and AC-03.2. \\ \hline
\end{longtable}
\endgroup
\par\medskip

\hypertarget{sec:b9-extension-points}{}
\addcontentsline{toc}{subsubsection}{B9. Extension Points}
\subsubsection*{B9. Extension Points}
\par\medskip

\begingroup
\small
\renewcommand{\baselinestretch}{1.0}\selectfont
\renewcommand{\arraystretch}{1.15}
\setlength{\tabcolsep}{6pt}
\begin{longtable}{|>{\raggedright\arraybackslash}p{0.39\textwidth}|>{\raggedright\arraybackslash}p{0.54\textwidth}|}
\hline
\textbf{Extension Point} & \textbf{Location in Flow of Events} \\ \hline
\endhead
Settings logout entry & Basic flow steps 1-2; alternative flow A1 \\ \hline
Global-navigation logout entry & Basic flow steps 1-2; alternative flow A2 \\ \hline
Client cleanup & Basic flow step 4 \\ \hline
Backend session invalidation & Basic flow step 5 \\ \hline
Redirect to Login & Basic flow step 6 \\ \hline
\end{longtable}
\endgroup
\par\medskip

\hypertarget{sec:b10-interface-and-data-notes}{}
\addcontentsline{toc}{subsubsection}{B10. Interface and Data Notes}
\subsubsection*{B10. Interface and Data Notes}
\par\medskip

\begin{itemize}
\item \textbf{Confirmed requirement --- screen/flow path:} \texttt{Home → Settings → Logout → Login}; \texttt{Home → Logout} from the screen flow; feature-flow logout follows the authenticated main pages.
\item \textbf{Confirmed requirement --- UI text:} The control is labeled \textbf{Logout} in \texttt{userstory.md} and \texttt{screen\_flow.md}, and \textbf{Log out} in \texttt{settings.md}.
\item \textbf{Confirmed requirement --- UI actions:} \texttt{LogoutButton} ends the current session and clears local/session authentication data; no confirmation dialog is specified for logout.
\item \textbf{Confirmed requirement --- data cleanup:} Authentication tokens and temporary user data in local and session storage must be cleared; backend session state must be invalidated.
\item \textbf{Data model note:} \texttt{settings.md} names \texttt{sessionId}, \texttt{startedAt}, \texttt{lastActivityAt}, and \texttt{expiresAt} as session fields; this is a settings information-model detail rather than an API field list.
\item \textbf{API endpoint:} Not specified in proposals; \texttt{API.md} has no logout endpoint.
\item \textbf{Statuses and error codes:} Not specified in proposals for logout itself; \texttt{API.md} generally defines \texttt{UNAUTHORIZED} 401 for protected endpoints, and no logout-specific error code appears in \texttt{response\_code.md}.
\item \textbf{Notification/privacy rules:} No logout-specific notification is specified; protected access after logout requires login according to NFR-SEC-1.
\item \textbf{Recommendation only:} \texttt{settings.md} recommends \texttt{/settings} and Account Settings grouping; the route is not defined in the proposal flow documents.
\end{itemize}
\par\medskip

\hypertarget{sec:b11-source-alignment-notes}{}
\addcontentsline{toc}{subsubsection}{B11. Source Alignment Notes}
\subsubsection*{B11. Source Alignment Notes}
\par\medskip

\begin{itemize}
\item \texttt{API.md} defines \texttt{POST /auth/login} but has no logout endpoint, while FR-03.1, FR-03.2, FR-03.3, AC-03.1, and AC-03.2 require logout, cleanup, and invalidation.
\item UC-02 Log In creates the authenticated session handled by UC-08 and already notes the missing logout endpoint; UC-08 does not repeat credential validation, CAPTCHA, or Google login behavior.
\item UC-01 Register Account can establish an initial authenticated session, but it does not define logout behavior; UC-08 applies to that session only after the user is logged in.
\item AC-03.3 and AC-03.4 are UI and navigation validation requirements; no exact screen mockup or logout message text is provided in the proposal sources.
\item No new I-xx, E-xx, or G-xx identifiers were created; the focused model uses only UC-08 and its primary actor.
\end{itemize}
\par\medskip

\pagebreak
\hypertarget{sec:uc-09-automatic-logout}{}
\addcontentsline{toc}{section}{UC-09 --- Automatic Logout}
\section*{UC-09 --- Automatic Logout}
\par\medskip

\textbf{Version:} \texttt{1.0}
\par\medskip

\textbf{Performed by:} [Nguyen Thanh An] | \textbf{Edited by:} [Nguyen Thanh An]
\par\medskip

\begingroup
\small
\renewcommand{\baselinestretch}{1.0}\selectfont
\renewcommand{\arraystretch}{1.15}
\setlength{\tabcolsep}{6pt}
\begin{longtable}{|>{\raggedright\arraybackslash}p{0.10\textwidth}|>{\raggedright\arraybackslash}p{0.83\textwidth}|}
\hline
\textbf{Field} & \textbf{Value} \\ \hline
\endhead
Use Case ID & UC-09 \\ \hline
Name & Automatic Logout \\ \hline
Primary Actor & user \\ \hline
Supporting Actors & None \\ \hline
Feature Group & Authentication Account \\ \hline
Priority & Not explicitly stated \\ \hline
Related User Stories & US-04 \\ \hline
Related Requirements & FR-04.1, FR-04.2, FR-04.3, FR-04.4; AC-04.1, AC-04.2, AC-04.3, AC-04.4; NFR-SEC-1, NFR-SEC-6, NFR-SEC-7 \\ \hline
Source Documents & \texttt{proposals/userstory.md}; \texttt{proposals/functional\_requirement.md}; \texttt{proposals/overview.md}; \texttt{proposals/feature\_flow.md}; \texttt{proposals/screen\_flow.md}; \texttt{proposals/API.md}; \texttt{proposals/response\_code.md}; \texttt{proposals/non\_functional\_requirement.md}; \texttt{proposals/settings.md} \\ \hline
\end{longtable}
\endgroup
\par\medskip

\hypertarget{sec:part-a-use-case-model}{}
\addcontentsline{toc}{subsection}{Part A --- Use Case Model}
\subsection*{Part A --- Use Case Model}
\par\medskip

\hypertarget{sec:a1-focused-use-case-diagram}{}
\addcontentsline{toc}{subsubsection}{A1. Focused Use Case Diagram}
\subsubsection*{A1. Focused Use Case Diagram}
\par\medskip

\begin{center}
\includegraphics[width=0.85\textwidth,height=0.75\textheight,keepaspectratio]{"figures/UC-09-automatic-logout_1.png"}
\end{center}
\par\medskip

\textbf{Diagram notation.} Person shape = human actor; cloud = supporting system; blue = selected use case; solid line = association; labelled solid arrow = generalization; dashed arrow = include or extend.
\par\medskip

\hypertarget{sec:a2-actor-register}{}
\addcontentsline{toc}{subsubsection}{A2. Actor Register}
\subsubsection*{A2. Actor Register}
\par\medskip

\begingroup
\small
\renewcommand{\baselinestretch}{1.0}\selectfont
\renewcommand{\arraystretch}{1.15}
\setlength{\tabcolsep}{6pt}
\begin{longtable}{|>{\raggedright\arraybackslash}p{0.09\textwidth}|>{\raggedright\arraybackslash}p{0.12\textwidth}|>{\raggedright\arraybackslash}p{0.70\textwidth}|}
\hline
\textbf{Actor} & \textbf{Type} & \textbf{Description} \\ \hline
\endhead
user & Primary human actor & A logged-in user whose inactive session is automatically terminated to protect the account when unattended. \\ \hline
\end{longtable}
\endgroup
\par\medskip

\hypertarget{sec:a3-relationships}{}
\addcontentsline{toc}{subsubsection}{A3. Relationships}
\subsubsection*{A3. Relationships}
\par\medskip

\hypertarget{sec:a3-1-associations}{}
\paragraph*{A3.1 Associations}
\par\medskip

\begingroup
\small
\renewcommand{\baselinestretch}{1.0}\selectfont
\renewcommand{\arraystretch}{1.15}
\setlength{\tabcolsep}{6pt}
\begin{longtable}{|>{\raggedright\arraybackslash}p{0.09\textwidth}|>{\raggedright\arraybackslash}p{0.15\textwidth}|>{\raggedright\arraybackslash}p{0.09\textwidth}|>{\raggedright\arraybackslash}p{0.58\textwidth}|}
\hline
\textbf{Actor} & \textbf{Use Case} & \textbf{Type} & \textbf{Description} \\ \hline
\endhead
user & UC-09 Automatic Logout & Primary & Receives automatic session expiration after inactivity and is required to log in again. \\ \hline
\end{longtable}
\endgroup
\par\medskip

\hypertarget{sec:a3-2-generalization-and-specialization}{}
\paragraph*{A3.2 Generalization and Specialization}
\par\medskip

\begingroup
\small
\renewcommand{\baselinestretch}{1.0}\selectfont
\renewcommand{\arraystretch}{1.15}
\setlength{\tabcolsep}{6pt}
\begin{longtable}{|>{\raggedright\arraybackslash}p{0.09\textwidth}|>{\raggedright\arraybackslash}p{0.09\textwidth}|>{\raggedright\arraybackslash}p{0.09\textwidth}|>{\raggedright\arraybackslash}p{0.64\textwidth}|}
\hline
\textbf{Child} & \textbf{Parent} & \textbf{Type} & \textbf{Inherited Behavior or Association} \\ \hline
\endhead
None & None & None & No actor or use-case generalization is specified in the proposals. \\ \hline
\end{longtable}
\endgroup
\par\medskip

\hypertarget{sec:a3-3-include}{}
\paragraph*{A3.3 Include}
\par\medskip

\begingroup
\small
\renewcommand{\baselinestretch}{1.0}\selectfont
\renewcommand{\arraystretch}{1.15}
\setlength{\tabcolsep}{6pt}
\begin{longtable}{|>{\raggedright\arraybackslash}p{0.16\textwidth}|>{\raggedright\arraybackslash}p{0.20\textwidth}|>{\raggedright\arraybackslash}p{0.56\textwidth}|}
\hline
\textbf{Base Use Case} & \textbf{Included Use Case} & \textbf{Meaning} \\ \hline
\endhead
None & None & No included use case is named in the proposals. \\ \hline
\end{longtable}
\endgroup
\par\medskip

\hypertarget{sec:a3-4-extend}{}
\paragraph*{A3.4 Extend}
\par\medskip

\begingroup
\small
\renewcommand{\baselinestretch}{1.0}\selectfont
\renewcommand{\arraystretch}{1.15}
\setlength{\tabcolsep}{6pt}
\begin{longtable}{|>{\raggedright\arraybackslash}p{0.19\textwidth}|>{\raggedright\arraybackslash}p{0.13\textwidth}|>{\raggedright\arraybackslash}p{0.10\textwidth}|>{\raggedright\arraybackslash}p{0.49\textwidth}|}
\hline
\textbf{Extending Use Case} & \textbf{Base Use Case} & \textbf{Condition} & \textbf{Added Behavior} \\ \hline
\endhead
None & None & None & No extending use case is named in the proposals. \\ \hline
\end{longtable}
\endgroup
\par\medskip

\hypertarget{sec:a4-model-notes}{}
\addcontentsline{toc}{subsubsection}{A4. Model Notes}
\subsubsection*{A4. Model Notes}
\par\medskip

\begin{enumerate}
\item UC-09 covers automatic termination after inactivity; manual logout is outside its scope and belongs to UC-08.
\item The 15-minute value comes from FR-04.1 and AC-04.1; NFR-SEC-7 states 30 minutes, so the conflict is recorded rather than silently resolved in the model.
\item Activity reset, protected-access denial, and expiration messaging are internal behavior because no separate include or extend use case is named.
\item No new I-xx, E-xx, or G-xx identifiers are introduced; the model contains only the primary actor and the selected base use case.
\end{enumerate}
\par\medskip

\hypertarget{sec:part-b-use-case-specification}{}
\addcontentsline{toc}{subsection}{Part B --- Use Case Specification}
\subsection*{Part B --- Use Case Specification}
\par\medskip

\hypertarget{sec:b1-brief-description}{}
\addcontentsline{toc}{subsubsection}{B1. Brief Description}
\subsubsection*{B1. Brief Description}
\par\medskip

The system automatically logs out an inactive user so an unattended account remains secure. After a successful login, no requests from the user for 15 consecutive minutes cause the session to terminate and the user to be logged out, as stated by FR-04.1 and AC-04.1.
\par\medskip

When the session has expired, protected pages and APIs deny access and require login again. The login redirect includes a message informing the user that the session expired; no exact message wording is provided in the proposals.
\par\medskip

\hypertarget{sec:b2-preconditions}{}
\addcontentsline{toc}{subsubsection}{B2. Preconditions}
\subsubsection*{B2. Preconditions}
\par\medskip

\begin{enumerate}
\item The user has successfully logged in and has an active session.
\item The inactivity timer is maintained by the frontend and backend as specified by FR-04.1 and FR-04.3.
\item The user does not send a qualifying request before the applicable timeout.
\end{enumerate}
\par\medskip

\hypertarget{sec:b3-postconditions}{}
\addcontentsline{toc}{subsubsection}{B3. Postconditions}
\subsubsection*{B3. Postconditions}
\par\medskip

\begin{enumerate}
\item \textbf{Success:} The inactive session is terminated and the user is logged out.
\item \textbf{Success:} The user is redirected to Login and informed that the session expired.
\item \textbf{Success:} Protected pages and APIs deny access using the expired session.
\item \textbf{Failure:} If a qualifying request is received before the timeout, the inactivity timer is reset and the session remains active.
\end{enumerate}
\par\medskip

\hypertarget{sec:b4-main-success-scenario-basic-flow}{}
\addcontentsline{toc}{subsubsection}{B4. Main Success Scenario (Basic Flow)}
\subsubsection*{B4. Main Success Scenario (Basic Flow)}
\par\medskip

\begin{enumerate}
\item The user completes login and receives an active session.
\item The system starts or maintains the inactivity timer for that session.
\item The system receives no requests from the user for 15 consecutive minutes, according to FR-04.1 and AC-04.1.
\item The system automatically terminates the session and logs the user out.
\item The user attempts to continue using the application or is redirected to Login.
\item The system displays a message informing the user that the session expired and requires login again.
\end{enumerate}
\par\medskip

\hypertarget{sec:b5-alternative-flows}{}
\addcontentsline{toc}{subsubsection}{B5. Alternative Flows}
\subsubsection*{B5. Alternative Flows}
\par\medskip

The story specifies no distinct alternative flow; the source-backed variant is activity before the timeout.
\par\medskip

\hypertarget{sec:a1-session-reset-on-activity}{}
\paragraph*{A1. Session reset on activity}
\par\medskip

At basic flow step 3, a new qualifying request is received before the timeout. The system resets the inactivity timer, keeps the session active, and resumes normal use of the protected page or API (FR-04.3, AC-04.3).
\par\medskip

\hypertarget{sec:b6-exception-flows}{}
\addcontentsline{toc}{subsubsection}{B6. Exception Flows}
\subsubsection*{B6. Exception Flows}
\par\medskip

\hypertarget{sec:e1-protected-page-access-after-session-expiration}{}
\paragraph*{E1. Protected page access after session expiration}
\par\medskip

The user attempts to open a protected page after expiry. The system denies access, redirects to Login, displays the session-expiration message, and requires the user to log in again (FR-04.2, FR-04.4, AC-04.2, AC-04.4).
\par\medskip

\hypertarget{sec:e2-protected-api-access-after-session-expiration}{}
\paragraph*{E2. Protected API access after session expiration}
\par\medskip

The user or frontend sends a protected API request after expiry. The system denies access and requires a new login; no automatic-renewal behavior is specified in the proposals (FR-04.2, AC-04.2).
\par\medskip

\hypertarget{sec:b7-business-rules}{}
\addcontentsline{toc}{subsubsection}{B7. Business Rules}
\subsubsection*{B7. Business Rules}
\par\medskip

\begin{enumerate}
\item After 15 consecutive minutes with no requests, the system automatically terminates the session and logs the user out (FR-04.1, AC-04.1).
\item NFR-SEC-7 states 30 minutes for inactive authentication sessions; this conflicts with FR-04.1 and AC-04.1 and is not silently overridden.
\item A protected page or API must deny access after session expiration and require login again (FR-04.2, AC-04.2).
\item A qualifying request received before the timeout resets the inactivity timer (FR-04.3, AC-04.3).
\item The login redirect after expiration must inform the user that the session expired (FR-04.4, AC-04.4).
\item Protected resources require authentication, and authorization checks must remain server-side (NFR-SEC-1, NFR-SEC-6).
\end{enumerate}
\par\medskip

\hypertarget{sec:b8-special-requirements}{}
\addcontentsline{toc}{subsubsection}{B8. Special Requirements}
\subsubsection*{B8. Special Requirements}
\par\medskip

\begingroup
\small
\renewcommand{\baselinestretch}{1.0}\selectfont
\renewcommand{\arraystretch}{1.15}
\setlength{\tabcolsep}{6pt}
\begin{longtable}{|>{\raggedright\arraybackslash}p{0.10\textwidth}|>{\raggedright\arraybackslash}p{0.83\textwidth}|}
\hline
\textbf{ID} & \textbf{Requirement} \\ \hline
\endhead
NFR-SEC-1 & Authentication is required before accessing protected user resources. \\ \hline
NFR-SEC-6 & Server-side authorization checks must prevent unauthorized access. \\ \hline
NFR-SEC-7 & Inactive authentication sessions must expire automatically after 30 minutes. \\ \hline
Design constraint & FR-04.1 and AC-04.1 require 15 minutes of inactivity; \texttt{settings.md} recommends using the stricter 15-minute value unless the product explicitly decides otherwise. \\ \hline
Design constraint & The expiration notice is required by FR-04.4 and AC-04.4, but its exact text is not specified in the proposals. \\ \hline
\end{longtable}
\endgroup
\par\medskip

\hypertarget{sec:b9-extension-points}{}
\addcontentsline{toc}{subsubsection}{B9. Extension Points}
\subsubsection*{B9. Extension Points}
\par\medskip

\begingroup
\small
\renewcommand{\baselinestretch}{1.0}\selectfont
\renewcommand{\arraystretch}{1.15}
\setlength{\tabcolsep}{6pt}
\begin{longtable}{|>{\raggedright\arraybackslash}p{0.47\textwidth}|>{\raggedright\arraybackslash}p{0.46\textwidth}|}
\hline
\textbf{Extension Point} & \textbf{Location in Flow of Events} \\ \hline
\endhead
Inactivity timer start or maintenance & Basic flow steps 1-2 \\ \hline
Timeout detection & Basic flow step 3 \\ \hline
Qualifying request & Alternative flow A1 \\ \hline
Protected-page denial & Exception flow E1 \\ \hline
Protected-API denial & Exception flow E2 \\ \hline
Expiration notice & Basic flow step 6; exception flow E1 \\ \hline
\end{longtable}
\endgroup
\par\medskip

\hypertarget{sec:b10-interface-and-data-notes}{}
\addcontentsline{toc}{subsubsection}{B10. Interface and Data Notes}
\subsubsection*{B10. Interface and Data Notes}
\par\medskip

\begin{itemize}
\item \textbf{Confirmed requirement --- screen/flow path:} Automatic logout returns the user to Login; \texttt{screen\_flow.md} shows \texttt{Home → Logout} and the Authentication area with \texttt{Login}.
\item \textbf{Confirmed requirement --- UI behavior:} The system informs the user that the session expired and requires login again; no exact quoted message text is specified.
\item \textbf{Confirmed requirement --- session data:} \texttt{settings.md} names \texttt{sessionId}, \texttt{startedAt}, \texttt{lastActivityAt}, and \texttt{expiresAt} for session summary and inactivity feedback.
\item \textbf{Confirmed requirement --- policy conflict handling:} \texttt{settings.md} recommends the stricter 15-minute value unless the product explicitly decides otherwise.
\item \textbf{API endpoint:} Not specified in proposals; \texttt{API.md} has no dedicated automatic-logout endpoint.
\item \textbf{Statuses and error codes:} Not specified in proposals for automatic logout itself; \texttt{API.md} generally defines \texttt{UNAUTHORIZED} 401, and no inactivity-specific error code appears in \texttt{response\_code.md}.
\item \textbf{Notification/privacy rules:} The only specified feedback is the session-expiration notice on redirect to Login; no email or notification-center message is specified.
\item \textbf{Recommendation only:} \texttt{/settings} and the Account Settings session summary are recommended UI groupings in \texttt{settings.md}, not proposal-defined routes.
\end{itemize}
\par\medskip

\hypertarget{sec:b11-source-alignment-notes}{}
\addcontentsline{toc}{subsubsection}{B11. Source Alignment Notes}
\subsubsection*{B11. Source Alignment Notes}
\par\medskip

\begin{itemize}
\item FR-04.1 and AC-04.1 require logout after 15 minutes of inactivity, while NFR-SEC-7 states 30 minutes; \texttt{settings.md} recommends the stricter 15-minute value unless the product explicitly decides otherwise.
\item UC-02 Log In already treats session expiry as a protected-access exception and lists NFR-SEC-7; UC-09 details the inactivity trigger, timer reset, denial, and notice without duplicating login credential behavior.
\item \texttt{API.md} has no dedicated automatic-logout endpoint; expired access is enforced through denial of protected pages and APIs rather than a documented route.
\item UC-01 Register Account can create an initial session, but it does not define inactivity policy; whether that session follows the same timeout is not specified in the proposals.
\item No new I-xx, E-xx, or G-xx identifiers were created; the focused model uses only UC-09 and its primary actor.
\end{itemize}
\par\medskip

\pagebreak
\hypertarget{sec:uc-10-recover-password}{}
\addcontentsline{toc}{section}{UC-10 --- Recover Password}
\section*{UC-10 --- Recover Password}
\par\medskip

\textbf{Version:} \texttt{1.0}
\par\medskip

\textbf{Performed by:} [Nguyen Thanh An] | \textbf{Edited by:} [Nguyen Thanh An]
\par\medskip

\begingroup
\small
\renewcommand{\baselinestretch}{1.0}\selectfont
\renewcommand{\arraystretch}{1.15}
\setlength{\tabcolsep}{6pt}
\begin{longtable}{|>{\raggedright\arraybackslash}p{0.10\textwidth}|>{\raggedright\arraybackslash}p{0.83\textwidth}|}
\hline
\textbf{Field} & \textbf{Value} \\ \hline
\endhead
Use Case ID & UC-10 \\ \hline
Name & Recover Password \\ \hline
Primary Actor & registered user \\ \hline
Supporting Actors & None \\ \hline
Feature Group & Authentication Account \\ \hline
Priority & Not explicitly stated \\ \hline
Related User Stories & US-05 \\ \hline
Related Requirements & FR-05.1, FR-05.2, FR-05.3, FR-05.4, FR-05.5; FR-07.1, FR-07.2, FR-07.3, FR-07.4, FR-07.5, FR-07.6; FR-01.9; AC-05.1, AC-05.2, AC-05.3, AC-05.4, AC-05.5, AC-05.6; AC-07.1, AC-07.2, AC-07.3, AC-07.4, AC-07.5, AC-07.6, AC-07.7, AC-07.8; NFR-SEC-1, NFR-SEC-2, NFR-SEC-3, NFR-SEC-8 \\ \hline
Source Documents & \texttt{proposals/userstory.md}; \texttt{proposals/functional\_requirement.md}; \texttt{proposals/overview.md}; \texttt{proposals/feature\_flow.md}; \texttt{proposals/screen\_flow.md}; \texttt{proposals/API.md}; \texttt{proposals/response\_code.md}; \texttt{proposals/non\_functional\_requirement.md}; \texttt{proposals/settings.md} \\ \hline
\end{longtable}
\endgroup
\par\medskip

\hypertarget{sec:part-a-use-case-model}{}
\addcontentsline{toc}{subsection}{Part A --- Use Case Model}
\subsection*{Part A --- Use Case Model}
\par\medskip

\hypertarget{sec:a1-focused-use-case-diagram}{}
\addcontentsline{toc}{subsubsection}{A1. Focused Use Case Diagram}
\subsubsection*{A1. Focused Use Case Diagram}
\par\medskip

\begin{center}
\includegraphics[width=0.85\textwidth,height=0.75\textheight,keepaspectratio]{"figures/UC-10-recover-password_1.png"}
\end{center}
\par\medskip

\textbf{Diagram notation.} Person shape = human actor; cloud = supporting system; blue = selected use case; solid line = association; labelled solid arrow = generalization; dashed arrow = include or extend.
\par\medskip

\hypertarget{sec:a2-actor-register}{}
\addcontentsline{toc}{subsubsection}{A2. Actor Register}
\subsubsection*{A2. Actor Register}
\par\medskip

\begingroup
\small
\renewcommand{\baselinestretch}{1.0}\selectfont
\renewcommand{\arraystretch}{1.15}
\setlength{\tabcolsep}{6pt}
\begin{longtable}{|>{\raggedright\arraybackslash}p{0.12\textwidth}|>{\raggedright\arraybackslash}p{0.15\textwidth}|>{\raggedright\arraybackslash}p{0.65\textwidth}|}
\hline
\textbf{Actor} & \textbf{Type} & \textbf{Description} \\ \hline
\endhead
registered user & Primary human actor & A user with an existing account who regains access after forgetting the password. \\ \hline
\end{longtable}
\endgroup
\par\medskip

\hypertarget{sec:a3-relationships}{}
\addcontentsline{toc}{subsubsection}{A3. Relationships}
\subsubsection*{A3. Relationships}
\par\medskip

\hypertarget{sec:a3-1-associations}{}
\paragraph*{A3.1 Associations}
\par\medskip

\begingroup
\small
\renewcommand{\baselinestretch}{1.0}\selectfont
\renewcommand{\arraystretch}{1.15}
\setlength{\tabcolsep}{6pt}
\begin{longtable}{|>{\raggedright\arraybackslash}p{0.09\textwidth}|>{\raggedright\arraybackslash}p{0.14\textwidth}|>{\raggedright\arraybackslash}p{0.09\textwidth}|>{\raggedright\arraybackslash}p{0.58\textwidth}|}
\hline
\textbf{Actor} & \textbf{Use Case} & \textbf{Type} & \textbf{Description} \\ \hline
\endhead
registered user & UC-10 Recover Password & Primary & Requests an OTP, sets a new password, and returns to Login after all sessions are invalidated. \\ \hline
\end{longtable}
\endgroup
\par\medskip

\hypertarget{sec:a3-2-generalization-and-specialization}{}
\paragraph*{A3.2 Generalization and Specialization}
\par\medskip

\begingroup
\small
\renewcommand{\baselinestretch}{1.0}\selectfont
\renewcommand{\arraystretch}{1.15}
\setlength{\tabcolsep}{6pt}
\begin{longtable}{|>{\raggedright\arraybackslash}p{0.09\textwidth}|>{\raggedright\arraybackslash}p{0.09\textwidth}|>{\raggedright\arraybackslash}p{0.09\textwidth}|>{\raggedright\arraybackslash}p{0.64\textwidth}|}
\hline
\textbf{Child} & \textbf{Parent} & \textbf{Type} & \textbf{Inherited Behavior or Association} \\ \hline
\endhead
None & None & None & No actor or use-case generalization is specified in the proposals. \\ \hline
\end{longtable}
\endgroup
\par\medskip

\hypertarget{sec:a3-3-include}{}
\paragraph*{A3.3 Include}
\par\medskip

\begingroup
\small
\renewcommand{\baselinestretch}{1.0}\selectfont
\renewcommand{\arraystretch}{1.15}
\setlength{\tabcolsep}{6pt}
\begin{longtable}{|>{\raggedright\arraybackslash}p{0.10\textwidth}|>{\raggedright\arraybackslash}p{0.12\textwidth}|>{\raggedright\arraybackslash}p{0.70\textwidth}|}
\hline
\textbf{Base Use Case} & \textbf{Included Use Case} & \textbf{Meaning} \\ \hline
\endhead
None & None & No included use case is named in the proposals; OTP rules are cited from US-07 without creating a new ID. \\ \hline
\end{longtable}
\endgroup
\par\medskip

\hypertarget{sec:a3-4-extend}{}
\paragraph*{A3.4 Extend}
\par\medskip

\begingroup
\small
\renewcommand{\baselinestretch}{1.0}\selectfont
\renewcommand{\arraystretch}{1.15}
\setlength{\tabcolsep}{6pt}
\begin{longtable}{|>{\raggedright\arraybackslash}p{0.19\textwidth}|>{\raggedright\arraybackslash}p{0.13\textwidth}|>{\raggedright\arraybackslash}p{0.10\textwidth}|>{\raggedright\arraybackslash}p{0.49\textwidth}|}
\hline
\textbf{Extending Use Case} & \textbf{Base Use Case} & \textbf{Condition} & \textbf{Added Behavior} \\ \hline
\endhead
None & None & None & No extending use case is named in the proposals. \\ \hline
\end{longtable}
\endgroup
\par\medskip

\hypertarget{sec:a4-model-notes}{}
\addcontentsline{toc}{subsubsection}{A4. Model Notes}
\subsubsection*{A4. Model Notes}
\par\medskip

\begin{enumerate}
\item UC-10 starts from Login via \textbf{Forget Button} or from Account Settings via \textbf{Recover Password}; it ends with a password reset and re-login.
\item OTP generation, delivery, expiry, verification, resend, and reuse prevention are governed by FR-07.1 to FR-07.6 and AC-07.1 to AC-07.8 but are not modeled as a new included use case.
\item Non-enumeration protection is internal backend behavior: the requester receives the same generic success message whether the account exists or not.
\item No new I-xx, E-xx, or G-xx identifiers are introduced; the model contains only the primary actor and the selected base use case.
\end{enumerate}
\par\medskip

\hypertarget{sec:part-b-use-case-specification}{}
\addcontentsline{toc}{subsection}{Part B --- Use Case Specification}
\subsection*{Part B --- Use Case Specification}
\par\medskip

\hypertarget{sec:b1-brief-description}{}
\addcontentsline{toc}{subsubsection}{B1. Brief Description}
\subsubsection*{B1. Brief Description}
\par\medskip

A registered user recovers or resets a forgotten password by submitting a registered email address or username and completing an email OTP challenge. The system protects against account enumeration by returning the same generic success message for existing and non-existent accounts.
\par\medskip

After a valid OTP, the user enters and confirms a new password that satisfies the password security policy. Completing the reset invalidates all active sessions across all devices and requires the user to log in again.
\par\medskip

\hypertarget{sec:b2-preconditions}{}
\addcontentsline{toc}{subsubsection}{B2. Preconditions}
\subsubsection*{B2. Preconditions}
\par\medskip

\begin{enumerate}
\item The user has forgotten the password but can access the registered email address.
\item The submitted identifier is a registered email address or username, or the request is treated as a non-existent account under the generic-response rule.
\item Email delivery and OTP verification are available.
\end{enumerate}
\par\medskip

\hypertarget{sec:b3-postconditions}{}
\addcontentsline{toc}{subsubsection}{B3. Postconditions}
\subsubsection*{B3. Postconditions}
\par\medskip

\begin{enumerate}
\item \textbf{Success:} The user's password is updated using a securely stored hash.
\item \textbf{Success:} All active sessions across all devices are invalidated, and the user must log in again.
\item \textbf{Success:} A non-existent account request produces the identical generic success message and creates no accessible session.
\item \textbf{Failure:} Invalid OTP, policy failure, or confirmation mismatch leaves the password unchanged.
\end{enumerate}
\par\medskip

\hypertarget{sec:b4-main-success-scenario-basic-flow}{}
\addcontentsline{toc}{subsubsection}{B4. Main Success Scenario (Basic Flow)}
\subsubsection*{B4. Main Success Scenario (Basic Flow)}
\par\medskip

\begin{enumerate}
\item The user selects \textbf{Forget Button} on Login or \textbf{Recover Password} in Account Settings.
\item The system displays the recovery form with \textbf{Username / Email} and \textbf{Send code}.
\item The user enters a registered email address or username and selects \textbf{Send code}.
\item The backend processes the request without revealing whether the account exists.
\item The system sends a time-limited OTP to the registered email address and shows the OTP verification screen.
\item The user enters the 6-digit OTP in the \textbf{6 Slot OTP} field and selects \textbf{Submit Button}.
\item The backend verifies that the OTP is valid, unexpired, and unused.
\item The system opens the \textbf{New Password} screen.
\item The user enters \textbf{New Password} and \textbf{Confirm New Password}.
\item The system verifies that the new password satisfies the password security policy and that the confirmation matches.
\item The system updates the stored password hash and invalidates all active sessions across all devices.
\item The system prompts the user to log in again.
\end{enumerate}
\par\medskip

\hypertarget{sec:b5-alternative-flows}{}
\addcontentsline{toc}{subsubsection}{B5. Alternative Flows}
\subsubsection*{B5. Alternative Flows}
\par\medskip

\hypertarget{sec:a1-recovery-request-for-a-non-existent-account}{}
\paragraph*{A1. Recovery request for a non-existent account}
\par\medskip

At basic flow step 4, the submitted identifier does not match an account. The system returns the identical generic success message required by FR-05.3 and AC-05.3; no OTP is delivered to a non-existent account.
\par\medskip

\hypertarget{sec:a2-resend-otp}{}
\paragraph*{A2. Resend OTP}
\par\medskip

At basic flow step 6, the user requests a new OTP. The system generates and sends a new 6-digit OTP, invalidates the previous OTP, and restarts the 5-minute expiry period (FR-07.6, AC-07.7).
\par\medskip

\hypertarget{sec:a3-change-password-entry-from-account-settings}{}
\paragraph*{A3. Change Password entry from Account Settings}
\par\medskip

At basic flow step 1, the user opens \textbf{Account} and selects \textbf{Change Password} rather than \textbf{Recover Password}. \texttt{settings.md} uses the same OTP-verified update sequence; the direct-change control is explicitly described as a recommendation derived from recovery policy.
\par\medskip

\hypertarget{sec:b6-exception-flows}{}
\addcontentsline{toc}{subsubsection}{B6. Exception Flows}
\subsubsection*{B6. Exception Flows}
\par\medskip

\hypertarget{sec:e1-otp-invalid-expired-or-already-used}{}
\paragraph*{E1. OTP invalid, expired, or already used}
\par\medskip

The system rejects the OTP, displays an appropriate error message, and does not open the password-reset form. The user may request a new OTP (FR-07.5, AC-07.4, AC-07.5, AC-07.8).
\par\medskip

\hypertarget{sec:e2-new-password-fails-the-security-policy}{}
\paragraph*{E2. New password fails the security policy}
\par\medskip

The system blocks submission and shows a validation error. The new password must contain at least 8 characters, including at least one lowercase letter, one uppercase letter, one number, and one special character (FR-05.4, FR-01.9, AC-05.4).
\par\medskip

\hypertarget{sec:e3-confirm-new-password-mismatch}{}
\paragraph*{E3. Confirm New Password mismatch}
\par\medskip

The system blocks submission until \textbf{Confirm New Password} matches \textbf{New Password} (AC-05.5).
\par\medskip

\hypertarget{sec:b7-business-rules}{}
\addcontentsline{toc}{subsubsection}{B7. Business Rules}
\subsubsection*{B7. Business Rules}
\par\medskip

\begin{enumerate}
\item The system accepts a registered email address or username (FR-05.1, AC-05.1).
\item The system sends a time-limited OTP to the registered email address (FR-05.2, AC-05.2).
\item A request for a non-existent account returns the identical generic success message to prevent user enumeration (FR-05.3, AC-05.3).
\item The new password must satisfy the password security policy: at least 8 characters with lowercase, uppercase, number, and special character (FR-05.4, FR-01.9, AC-05.4).
\item \textbf{Confirm New Password} must match \textbf{New Password} (AC-05.5).
\item Completing a password reset logs out all active sessions across all devices (FR-05.5, AC-05.6).
\item The OTP is 6 digits, expires 5 minutes after generation, is accepted only while valid and unused, and is invalidated on resend (FR-07.1, FR-07.2, FR-07.3, FR-07.4, FR-07.6; AC-07.1, AC-07.2, AC-07.3, AC-07.6, AC-07.7, AC-07.8).
\end{enumerate}
\par\medskip

\hypertarget{sec:b8-special-requirements}{}
\addcontentsline{toc}{subsubsection}{B8. Special Requirements}
\subsubsection*{B8. Special Requirements}
\par\medskip

\begingroup
\small
\renewcommand{\baselinestretch}{1.0}\selectfont
\renewcommand{\arraystretch}{1.15}
\setlength{\tabcolsep}{6pt}
\begin{longtable}{|>{\raggedright\arraybackslash}p{0.11\textwidth}|>{\raggedright\arraybackslash}p{0.82\textwidth}|}
\hline
\textbf{ID} & \textbf{Requirement} \\ \hline
\endhead
NFR-SEC-1 & Authentication is required before accessing protected user resources. \\ \hline
NFR-SEC-2 & User credentials must be stored using one-way password hashing with a secure hashing algorithm. \\ \hline
NFR-SEC-3 & Sensitive communication must be encrypted using TLS 1.2 or higher. \\ \hline
NFR-SEC-8 & The system must protect against common web security vulnerabilities, including SQL injection, XSS, and CSRF. \\ \hline
Design constraint & The OTP is a 6-digit numeric code, expires after 5 minutes, is single-use, and is invalidated when a new OTP is generated. \\ \hline
Design constraint & Password hashes must never be displayed; the generic recovery response must not distinguish existing from non-existent accounts. \\ \hline
\end{longtable}
\endgroup
\par\medskip

\hypertarget{sec:b9-extension-points}{}
\addcontentsline{toc}{subsubsection}{B9. Extension Points}
\subsubsection*{B9. Extension Points}
\par\medskip

\begingroup
\small
\renewcommand{\baselinestretch}{1.0}\selectfont
\renewcommand{\arraystretch}{1.15}
\setlength{\tabcolsep}{6pt}
\begin{longtable}{|>{\raggedright\arraybackslash}p{0.35\textwidth}|>{\raggedright\arraybackslash}p{0.58\textwidth}|}
\hline
\textbf{Extension Point} & \textbf{Location in Flow of Events} \\ \hline
\endhead
Login or Account Settings entry & Basic flow step 1; alternative flow A3 \\ \hline
Generic non-enumeration response & Basic flow step 4; alternative flow A1 \\ \hline
OTP verification & Basic flow steps 5-7; alternative flow A2; exception E1 \\ \hline
New password entry & Basic flow steps 8-10; exceptions E2-E3 \\ \hline
Session invalidation and re-login & Basic flow steps 11-12 \\ \hline
\end{longtable}
\endgroup
\par\medskip

\hypertarget{sec:b10-interface-and-data-notes}{}
\addcontentsline{toc}{subsubsection}{B10. Interface and Data Notes}
\subsubsection*{B10. Interface and Data Notes}
\par\medskip

\begin{itemize}
\item \textbf{Confirmed requirement --- screen/flow path:} \texttt{Login → Recover → OTP → New Password → Login} after reset; \texttt{Login → Recover} is shown as \texttt{Forget Button} in \texttt{userstory.md}; \texttt{screen\_flow.md} shows \texttt{Login .-$>$ Recover}; \texttt{feature\_flow.md} shows \texttt{Authentication → Recover Password}.
\item \textbf{Confirmed requirement --- recovery UI:} \textbf{Username / Email}, \textbf{Send code}, \textbf{6 Slot OTP}, \textbf{Resend code}, and \textbf{Submit Button} are named in \texttt{userstory.md}.
\item \textbf{Confirmed requirement --- new-password UI:} \textbf{New Password} and \textbf{Confirm New Password} are named in \texttt{userstory.md}; \texttt{settings.md} names \texttt{NewPasswordInput}, \texttt{ConfirmPasswordInput}, and \texttt{PasswordVisibilityToggle}.
\item \textbf{Confirmed requirement --- OTP notice:} The generic OTP UI in \texttt{userstory.md} shows \textbf{Sent OTP to email ab}*\textbf{@gmail.com}; recovery sends the code only to the registered email address.
\item \textbf{Confirmed requirement --- data fields:} Submitted identifier, OTP, new password, confirmation, password hash, and session invalidation across all devices.
\item \textbf{API endpoint:} Not specified in proposals; \texttt{API.md} has no password-recovery endpoint.
\item \textbf{Statuses and error codes:} Not specified in proposals for recovery; \texttt{response\_code.md} lists \texttt{INVALID\_EMAIL} 45002, \texttt{INVALID\_USERNAME} 45003, \texttt{INVALID\_PASSWORD} 45004, backend \texttt{UNKNOWN ERROR} 2002, and \texttt{UNKNOWN DATABASE ERROR} 45000, but no OTP-specific code.
\item \textbf{Notification/privacy rules:} OTP is delivered to the registered email; non-existent accounts receive the identical generic success message; password hashes are never displayed.
\item \textbf{Recommendation only:} \texttt{/settings/account}, \texttt{PasswordChangeButton}, and direct \textbf{Change Password} are Settings recommendations, not proposal-defined API routes.
\end{itemize}
\par\medskip

\hypertarget{sec:b11-source-alignment-notes}{}
\addcontentsline{toc}{subsubsection}{B11. Source Alignment Notes}
\subsubsection*{B11. Source Alignment Notes}
\par\medskip

\begin{itemize}
\item \texttt{API.md} has no password-recovery endpoint, while FR-05.1 to FR-05.5 and AC-05.1 to AC-05.6 require recovery, OTP delivery, validation, and session invalidation.
\item UC-01 Register Account already defines the password policy and OTP rules and notes that OTP submission and resend endpoints are absent; UC-10 reuses those exact FR/AC IDs without creating a new OTP use case.
\item UC-02 Log In provides the \textbf{Forget Button} entry into Recover; UC-10 begins after that navigation and does not repeat login credential or CAPTCHA behavior.
\item \texttt{userstory.md} draws \texttt{New Password → Submit Button → Home}, while AC-05.6, FR-05.5, and \texttt{settings.md} 7.4 require all sessions to be invalidated and the user to log in again; this specification follows the explicit AC/FR/settings behavior.
\item No new I-xx, E-xx, or G-xx identifiers were created; the focused model uses only UC-10 and its primary actor.
\end{itemize}
\par\medskip

\par\medskip
\hypertarget{sec:ui-prototype-specimens}{}
\addcontentsline{toc}{subsection}{UI Prototype Specimens}
\subsection*{UI Prototype Specimens}
\par\medskip
\begin{figure}[H]
\centering
\includegraphics[width=0.75\textwidth,height=0.42\textheight,keepaspectratio]{"../../prototype/otp_desktop.png"}
\par\vspace{2pt}
{\footnotesize \textbf{UI Prototype:} OTP - Desktop}
\end{figure}
\par\smallskip
\begin{figure}[H]
\centering
\includegraphics[width=0.75\textwidth,height=0.42\textheight,keepaspectratio]{"../../prototype/recover_desktop.png"}
\par\vspace{2pt}
{\footnotesize \textbf{UI Prototype:} Recover - Desktop}
\end{figure}
\par\smallskip
\begin{figure}[H]
\centering
\includegraphics[width=0.75\textwidth,height=0.42\textheight,keepaspectratio]{"../../prototype/recovery_confirmation_desktop.png"}
\par\vspace{2pt}
{\footnotesize \textbf{UI Prototype:} Recovery confirmation - Desktop}
\end{figure}
\par\smallskip
\begin{figure}[H]
\centering
\includegraphics[width=0.75\textwidth,height=0.42\textheight,keepaspectratio]{"../../prototype/set_new_password_desktop.png"}
\par\vspace{2pt}
{\footnotesize \textbf{UI Prototype:} Set new password - Desktop}
\end{figure}
\par\smallskip

\pagebreak
\hypertarget{sec:uc-11-delete-account}{}
\addcontentsline{toc}{section}{UC-11 --- Delete Account}
\section*{UC-11 --- Delete Account}
\par\medskip

\textbf{Version:} \texttt{1.0}
\par\medskip

\textbf{Performed by:} [Nguyen Thanh An] | \textbf{Edited by:} [Nguyen Thanh An]
\par\medskip

\begingroup
\small
\renewcommand{\baselinestretch}{1.0}\selectfont
\renewcommand{\arraystretch}{1.15}
\setlength{\tabcolsep}{6pt}
\begin{longtable}{|>{\raggedright\arraybackslash}p{0.10\textwidth}|>{\raggedright\arraybackslash}p{0.83\textwidth}|}
\hline
\textbf{Field} & \textbf{Value} \\ \hline
\endhead
Use Case ID & UC-11 \\ \hline
Name & Delete Account \\ \hline
Primary Actor & registered user \\ \hline
Supporting Actors & None \\ \hline
Feature Group & Authentication Account \\ \hline
Priority & Not explicitly stated \\ \hline
Related User Stories & US-06 \\ \hline
Related Requirements & FR-06.1, FR-06.2, FR-06.3, FR-06.4, FR-06.5, FR-06.6; FR-07.1, FR-07.2, FR-07.3, FR-07.4, FR-07.5, FR-07.6; AC-06.1, AC-06.2, AC-06.3, AC-06.4, AC-06.5, AC-06.6; AC-07.1, AC-07.2, AC-07.3, AC-07.4, AC-07.5, AC-07.6, AC-07.7, AC-07.8; NFR-SEC-1, NFR-SEC-2, NFR-SEC-3, NFR-SEC-6, NFR-SEC-8 \\ \hline
Source Documents & \texttt{proposals/userstory.md}; \texttt{proposals/functional\_requirement.md}; \texttt{proposals/overview.md}; \texttt{proposals/feature\_flow.md}; \texttt{proposals/screen\_flow.md}; \texttt{proposals/API.md}; \texttt{proposals/response\_code.md}; \texttt{proposals/non\_functional\_requirement.md}; \texttt{proposals/settings.md} \\ \hline
\end{longtable}
\endgroup
\par\medskip

\hypertarget{sec:part-a-use-case-model}{}
\addcontentsline{toc}{subsection}{Part A --- Use Case Model}
\subsection*{Part A --- Use Case Model}
\par\medskip

\hypertarget{sec:a1-focused-use-case-diagram}{}
\addcontentsline{toc}{subsubsection}{A1. Focused Use Case Diagram}
\subsubsection*{A1. Focused Use Case Diagram}
\par\medskip

\begin{center}
\includegraphics[width=0.85\textwidth,height=0.75\textheight,keepaspectratio]{"figures/UC-11-delete-account_1.png"}
\end{center}
\par\medskip

\textbf{Diagram notation.} Person shape = human actor; cloud = supporting system; blue = selected use case; solid line = association; labelled solid arrow = generalization; dashed arrow = include or extend.
\par\medskip

\hypertarget{sec:a2-actor-register}{}
\addcontentsline{toc}{subsubsection}{A2. Actor Register}
\subsubsection*{A2. Actor Register}
\par\medskip

\begingroup
\small
\renewcommand{\baselinestretch}{1.0}\selectfont
\renewcommand{\arraystretch}{1.15}
\setlength{\tabcolsep}{6pt}
\begin{longtable}{|>{\raggedright\arraybackslash}p{0.10\textwidth}|>{\raggedright\arraybackslash}p{0.11\textwidth}|>{\raggedright\arraybackslash}p{0.71\textwidth}|}
\hline
\textbf{Actor} & \textbf{Type} & \textbf{Description} \\ \hline
\endhead
registered user & Primary human actor & A user who requests permanent deletion of an account so personal information, posts, and data are erased from the platform. \\ \hline
\end{longtable}
\endgroup
\par\medskip

\hypertarget{sec:a3-relationships}{}
\addcontentsline{toc}{subsubsection}{A3. Relationships}
\subsubsection*{A3. Relationships}
\par\medskip

\hypertarget{sec:a3-1-associations}{}
\paragraph*{A3.1 Associations}
\par\medskip

\begingroup
\small
\renewcommand{\baselinestretch}{1.0}\selectfont
\renewcommand{\arraystretch}{1.15}
\setlength{\tabcolsep}{6pt}
\begin{longtable}{|>{\raggedright\arraybackslash}p{0.10\textwidth}|>{\raggedright\arraybackslash}p{0.13\textwidth}|>{\raggedright\arraybackslash}p{0.09\textwidth}|>{\raggedright\arraybackslash}p{0.58\textwidth}|}
\hline
\textbf{Actor} & \textbf{Use Case} & \textbf{Type} & \textbf{Description} \\ \hline
\endhead
registered user & UC-11 Delete Account & Primary & Confirms deletion, completes email OTP re-authentication, and loses all account access. \\ \hline
\end{longtable}
\endgroup
\par\medskip

\hypertarget{sec:a3-2-generalization-and-specialization}{}
\paragraph*{A3.2 Generalization and Specialization}
\par\medskip

\begingroup
\small
\renewcommand{\baselinestretch}{1.0}\selectfont
\renewcommand{\arraystretch}{1.15}
\setlength{\tabcolsep}{6pt}
\begin{longtable}{|>{\raggedright\arraybackslash}p{0.09\textwidth}|>{\raggedright\arraybackslash}p{0.09\textwidth}|>{\raggedright\arraybackslash}p{0.09\textwidth}|>{\raggedright\arraybackslash}p{0.64\textwidth}|}
\hline
\textbf{Child} & \textbf{Parent} & \textbf{Type} & \textbf{Inherited Behavior or Association} \\ \hline
\endhead
None & None & None & No actor or use-case generalization is specified in the proposals. \\ \hline
\end{longtable}
\endgroup
\par\medskip

\hypertarget{sec:a3-3-include}{}
\paragraph*{A3.3 Include}
\par\medskip

\begingroup
\small
\renewcommand{\baselinestretch}{1.0}\selectfont
\renewcommand{\arraystretch}{1.15}
\setlength{\tabcolsep}{6pt}
\begin{longtable}{|>{\raggedright\arraybackslash}p{0.10\textwidth}|>{\raggedright\arraybackslash}p{0.11\textwidth}|>{\raggedright\arraybackslash}p{0.71\textwidth}|}
\hline
\textbf{Base Use Case} & \textbf{Included Use Case} & \textbf{Meaning} \\ \hline
\endhead
None & None & No included use case is named in the proposals; email OTP rules are cited from US-07 without creating a new ID. \\ \hline
\end{longtable}
\endgroup
\par\medskip

\hypertarget{sec:a3-4-extend}{}
\paragraph*{A3.4 Extend}
\par\medskip

\begingroup
\small
\renewcommand{\baselinestretch}{1.0}\selectfont
\renewcommand{\arraystretch}{1.15}
\setlength{\tabcolsep}{6pt}
\begin{longtable}{|>{\raggedright\arraybackslash}p{0.19\textwidth}|>{\raggedright\arraybackslash}p{0.13\textwidth}|>{\raggedright\arraybackslash}p{0.10\textwidth}|>{\raggedright\arraybackslash}p{0.49\textwidth}|}
\hline
\textbf{Extending Use Case} & \textbf{Base Use Case} & \textbf{Condition} & \textbf{Added Behavior} \\ \hline
\endhead
None & None & None & No extending use case is named in the proposals. \\ \hline
\end{longtable}
\endgroup
\par\medskip

\hypertarget{sec:a4-model-notes}{}
\addcontentsline{toc}{subsubsection}{A4. Model Notes}
\subsubsection*{A4. Model Notes}
\par\medskip

\begin{enumerate}
\item UC-11 covers permanent account deletion only; deletion of a published resource is outside its scope.
\item Confirmation, cancellation, email OTP re-authentication, final deletion, session invalidation, and access rejection are internal behavior of UC-11 because no separate include or extend use case is named.
\item The OTP flow follows US-07 rules and the Settings re-authentication flow, but no new OTP use-case identifier is introduced.
\item No new I-xx, E-xx, or G-xx identifiers are introduced; the model contains only the primary actor and the selected base use case.
\end{enumerate}
\par\medskip

\hypertarget{sec:part-b-use-case-specification}{}
\addcontentsline{toc}{subsection}{Part B --- Use Case Specification}
\subsection*{Part B --- Use Case Specification}
\par\medskip

\hypertarget{sec:b1-brief-description}{}
\addcontentsline{toc}{subsubsection}{B1. Brief Description}
\subsubsection*{B1. Brief Description}
\par\medskip

A registered user permanently deletes an account so personal information, posts, and data are erased from the platform. The system first displays a confirmation dialog that explains the consequences, allows cancellation, and requires email OTP re-authentication before the destructive action can proceed.
\par\medskip

After a valid OTP and final confirmation, the account is deleted, all active sessions across all devices are logged out, and future access to the deleted account is rejected. The Settings flow names this final action \textbf{Permanently Delete}.
\par\medskip

\hypertarget{sec:b2-preconditions}{}
\addcontentsline{toc}{subsubsection}{B2. Preconditions}
\subsubsection*{B2. Preconditions}
\par\medskip

\begin{enumerate}
\item The user has a registered account and an active authenticated session.
\item The user can access the registered email address for OTP re-authentication.
\item The user opens the Account area and can reach \textbf{Delete Account}.
\end{enumerate}
\par\medskip

\hypertarget{sec:b3-postconditions}{}
\addcontentsline{toc}{subsubsection}{B3. Postconditions}
\subsubsection*{B3. Postconditions}
\par\medskip

\begin{enumerate}
\item \textbf{Success:} The account is permanently deleted, and personal information, posts, and data are erased from the platform.
\item \textbf{Success:} All active sessions across all devices are invalidated.
\item \textbf{Success:} Any attempt to access the deleted account is rejected.
\item \textbf{Success:} The user is returned to Login and can no longer authenticate with the deleted account.
\item \textbf{Failure:} If the user cancels or provides an invalid OTP, the account remains active.
\end{enumerate}
\par\medskip

\hypertarget{sec:b4-main-success-scenario-basic-flow}{}
\addcontentsline{toc}{subsubsection}{B4. Main Success Scenario (Basic Flow)}
\subsubsection*{B4. Main Success Scenario (Basic Flow)}
\par\medskip

\begin{enumerate}
\item The user opens \textbf{Account} in Settings and selects \textbf{Delete Account}.
\item The system displays \texttt{DeleteAccountConfirmDialog}, explaining the consequences of deletion.
\item The user continues past the confirmation dialog.
\item The system requires re-authentication through email OTP before permanently deleting the account.
\item The system sends a 6-digit OTP to the registered email address and displays \texttt{ReauthOtpDialog} with the OTP input, resend control, and expiry hint.
\item The user submits a valid, unexpired OTP and selects \textbf{Permanently Delete}.
\item The system verifies that the OTP is valid, unexpired, and unused.
\item The system permanently deletes the account and erases personal information, posts, and data from the platform.
\item The system invalidates all active sessions across all devices and rejects future access to the deleted account.
\item The system returns the user to Login.
\end{enumerate}
\par\medskip

\hypertarget{sec:b5-alternative-flows}{}
\addcontentsline{toc}{subsubsection}{B5. Alternative Flows}
\subsubsection*{B5. Alternative Flows}
\par\medskip

\hypertarget{sec:a1-cancel-deletion}{}
\paragraph*{A1. Cancel deletion}
\par\medskip

At basic flow step 2 or 3, the user cancels deletion. The system closes the deletion flow and the account remains active (FR-06.3, AC-06.2).
\par\medskip

\hypertarget{sec:a2-resend-otp}{}
\paragraph*{A2. Resend OTP}
\par\medskip

At basic flow step 6, the user requests a new OTP. The system generates and sends a new 6-digit OTP, invalidates the previous OTP, and restarts the 5-minute expiry period (FR-07.6, AC-07.7).
\par\medskip

\hypertarget{sec:b6-exception-flows}{}
\addcontentsline{toc}{subsubsection}{B6. Exception Flows}
\subsubsection*{B6. Exception Flows}
\par\medskip

\hypertarget{sec:e1-otp-invalid-expired-or-already-used}{}
\paragraph*{E1. OTP invalid, expired, or already used}
\par\medskip

The system rejects the OTP, displays an appropriate message, and does not permanently delete the account. The user may request a new OTP (FR-07.5, AC-07.4, AC-07.5, AC-07.8).
\par\medskip

\hypertarget{sec:e2-access-to-a-deleted-account}{}
\paragraph*{E2. Access to a deleted account}
\par\medskip

If the user or another client attempts to access the deleted account, the system rejects the request and keeps all sessions invalidated (FR-06.5, AC-06.5, FR-06.6, AC-06.6).
\par\medskip

\hypertarget{sec:b7-business-rules}{}
\addcontentsline{toc}{subsubsection}{B7. Business Rules}
\subsubsection*{B7. Business Rules}
\par\medskip

\begin{enumerate}
\item The system allows a user to request permanent deletion of the account (FR-06.1).
\item Selecting \textbf{Delete Account} displays a confirmation dialog explaining the consequences of deletion (FR-06.2, AC-06.1).
\item The user can cancel deletion, and the account must remain active (FR-06.3, AC-06.2).
\item When the user confirms deletion, the system proceeds with the account deletion process (AC-06.3).
\item The system requires re-authentication through email OTP before permanently deleting the account (FR-06.4, AC-06.4).
\item The OTP is 6 digits, expires 5 minutes after generation, is accepted only while valid and unused, and is invalidated on resend (FR-07.1, FR-07.2, FR-07.3, FR-07.4, FR-07.6; AC-07.1, AC-07.2, AC-07.3, AC-07.6, AC-07.7, AC-07.8).
\item Deleting an account logs out all active sessions across all devices and rejects access to the deleted account (FR-06.5, FR-06.6, AC-06.5, AC-06.6).
\item The final destructive action is \textbf{Permanently Delete} and is available only after confirmation and a valid OTP, as specified in the Settings deletion flow.
\end{enumerate}
\par\medskip

\hypertarget{sec:b8-special-requirements}{}
\addcontentsline{toc}{subsubsection}{B8. Special Requirements}
\subsubsection*{B8. Special Requirements}
\par\medskip

\begingroup
\small
\renewcommand{\baselinestretch}{1.0}\selectfont
\renewcommand{\arraystretch}{1.15}
\setlength{\tabcolsep}{6pt}
\begin{longtable}{|>{\raggedright\arraybackslash}p{0.10\textwidth}|>{\raggedright\arraybackslash}p{0.83\textwidth}|}
\hline
\textbf{ID} & \textbf{Requirement} \\ \hline
\endhead
NFR-SEC-1 & Authentication is required before accessing protected user resources. \\ \hline
NFR-SEC-2 & User credentials must be stored using one-way password hashing with a secure hashing algorithm. \\ \hline
NFR-SEC-3 & Sensitive communication must be encrypted using TLS 1.2 or higher. \\ \hline
NFR-SEC-6 & Server-side authorization checks must prevent unauthorized access. \\ \hline
NFR-SEC-8 & The system must protect against common web security vulnerabilities, including SQL injection, XSS, and CSRF. \\ \hline
Design constraint & OTP re-authentication is mandatory before permanent deletion; the OTP is 6 digits, expires after 5 minutes, is single-use, and is invalidated when a new OTP is generated. \\ \hline
Design constraint & Deleted-account access must be rejected and all sessions invalidated after deletion. \\ \hline
\end{longtable}
\endgroup
\par\medskip

\hypertarget{sec:b9-extension-points}{}
\addcontentsline{toc}{subsubsection}{B9. Extension Points}
\subsubsection*{B9. Extension Points}
\par\medskip

\begingroup
\small
\renewcommand{\baselinestretch}{1.0}\selectfont
\renewcommand{\arraystretch}{1.15}
\setlength{\tabcolsep}{6pt}
\begin{longtable}{|>{\raggedright\arraybackslash}p{0.34\textwidth}|>{\raggedright\arraybackslash}p{0.59\textwidth}|}
\hline
\textbf{Extension Point} & \textbf{Location in Flow of Events} \\ \hline
\endhead
Account Settings deletion entry & Basic flow step 1 \\ \hline
Confirmation and cancellation & Basic flow steps 2-3; alternative flow A1 \\ \hline
Email OTP re-authentication & Basic flow steps 4-7; alternative flow A2; exception E1 \\ \hline
Final permanent deletion & Basic flow steps 8-10 \\ \hline
Deleted-account access rejection & Exception flow E2 \\ \hline
\end{longtable}
\endgroup
\par\medskip

\hypertarget{sec:b10-interface-and-data-notes}{}
\addcontentsline{toc}{subsubsection}{B10. Interface and Data Notes}
\subsubsection*{B10. Interface and Data Notes}
\par\medskip

\begin{itemize}
\item \textbf{Confirmed requirement --- screen/flow path:} \texttt{Home → Profile → Re-auth → Delete → Login} in \texttt{userstory.md}; \texttt{settings.md} 7.5 gives \texttt{Account → Delete Account → confirmation → email OTP → Permanently Delete → session rejection}.
\item \textbf{Confirmed requirement --- UI text:} \textbf{Delete Account}, \textbf{Permanently Delete}, \texttt{DeleteAccountButton}, \texttt{DeleteAccountConfirmDialog}, \texttt{ReauthOtpDialog}, and \texttt{DeleteAccountFinalConfirmButton} are named in \texttt{settings.md}.
\item \textbf{Confirmed requirement --- OTP UI:} \texttt{ReauthOtpDialog} shows a 6-digit OTP input, resend control, and expiry hint; the generic OTP notice in \texttt{userstory.md} reads \textbf{Sent OTP to email ab}*\textbf{@gmail.com}.
\item \textbf{Confirmed requirement --- confirmation:} \texttt{DeleteAccountConfirmDialog} explains the consequences of deletion and provides cancel/confirm actions; no exact consequence-paragraph text is specified.
\item \textbf{Confirmed requirement --- account data:} \texttt{settings.md} names \texttt{userId}, \texttt{email}, \texttt{emailVerified}, \texttt{status}, and \texttt{linkedProviders}; \texttt{API.md} defines status values \texttt{0 (inactive)}, \texttt{1(preprofile)}, \texttt{2(active)}, and \texttt{3(suspend)}.
\item \textbf{Confirmed requirement --- deletion outcome:} Personal information, posts, and data are erased from the platform; future access is rejected; all sessions are invalidated.
\item \textbf{API endpoint:} Not specified in proposals; \texttt{API.md} has no account-deletion endpoint.
\item \textbf{Statuses and error codes:} Not specified in proposals for deletion; \texttt{API.md} generally defines \texttt{UNAUTHORIZED} 401 and \texttt{FORBIDDEN} 403, and \texttt{response\_code.md} lists \texttt{FORBIDDEN} 48010, backend \texttt{UNKNOWN ERROR} 2002, and \texttt{UNKNOWN DATABASE ERROR} 45000, but no deletion-specific code.
\item \textbf{Recommendation only:} \texttt{/settings/account} is a recommended route in \texttt{settings.md}; the Settings flow states personal data is deleted or anonymized according to platform policy, while the story requires erasure from the platform.
\item \textbf{Notification/privacy rules:} OTP is delivered to the registered email; no notification to other users is specified; deleted-account access must be rejected.
\end{itemize}
\par\medskip

\hypertarget{sec:b11-source-alignment-notes}{}
\addcontentsline{toc}{subsubsection}{B11. Source Alignment Notes}
\subsubsection*{B11. Source Alignment Notes}
\par\medskip

\begin{itemize}
\item \texttt{API.md} has no account-deletion endpoint and no OTP submission or resend endpoint, while FR-06.1 to FR-06.6 and AC-06.1 to AC-06.6 require deletion, OTP re-authentication, session invalidation, and access rejection.
\item \texttt{settings.md} 7.5 supplies the email OTP re-authentication and permanent deletion flow: select \textbf{Delete Account}, review \texttt{DeleteAccountConfirmDialog}, cancel or continue, submit a valid email OTP, select \textbf{Permanently Delete}, then invalidate sessions and reject access.
\item UC-01 Register Account defines OTP validation for email verification; UC-11 applies the same US-07 FR/AC rules to deletion re-authentication without reusing or creating a numbered OTP use-case ID.
\item UC-02 Log In ends by creating the sessions that UC-11 invalidates; UC-11 returns the user to Login and does not repeat credential validation, CAPTCHA, or Google login behavior.
\item No new I-xx, E-xx, or G-xx identifiers were created; the focused model uses only UC-11 and its primary actor.
\end{itemize}
\par\medskip

\par\medskip
\hypertarget{sec:ui-prototype-specimens}{}
\addcontentsline{toc}{subsection}{UI Prototype Specimens}
\subsection*{UI Prototype Specimens}
\par\medskip
\begin{figure}[H]
\centering
\includegraphics[width=0.75\textwidth,height=0.42\textheight,keepaspectratio]{"../../prototype/settings_03_settings_account.png"}
\par\vspace{2pt}
{\footnotesize \textbf{UI Prototype:} Settings: 03 Settings - Account}
\end{figure}
\par\smallskip
\begin{figure}[H]
\centering
\includegraphics[width=0.75\textwidth,height=0.42\textheight,keepaspectratio]{"../../prototype/settings_delete_account_desktop_otp_error.png"}
\par\vspace{2pt}
{\footnotesize \textbf{UI Prototype:} Settings: Delete Account - Desktop OTP error}
\end{figure}
\par\smallskip
\begin{figure}[H]
\centering
\includegraphics[width=0.75\textwidth,height=0.42\textheight,keepaspectratio]{"../../prototype/settings_delete_account_desktop_otp.png"}
\par\vspace{2pt}
{\footnotesize \textbf{UI Prototype:} Settings: Delete Account - Desktop OTP}
\end{figure}
\par\smallskip
\begin{figure}[H]
\centering
\includegraphics[width=0.75\textwidth,height=0.42\textheight,keepaspectratio]{"../../prototype/settings_delete_account_desktop_success.png"}
\par\vspace{2pt}
{\footnotesize \textbf{UI Prototype:} Settings: Delete Account - Desktop success}
\end{figure}
\par\smallskip

\end{document}
